\section{Harmonic sampling and stopped-interface geometry}
\label{sec:geometry}

We use lattice units in this section. A strand observation specifies the
orientation of finitely many pieces of zero interface, hence only the signs
of adjacent heights. Its conditional law is uniform on the corresponding
convex height polytope. All assertions below concern admissible observations,
that is, observations of positive probability. Constants are uniform in
the number of lattice edges in an observed strand. The number of distinct
macroscopic strands in a comparison is bounded. The exterior heights have
absolute value at most \(1/2\).

The geometric constructions and the random-walk estimates used below are
those of \cite[Sections 2.3 and 3.3--3.7]{SSdiscrete}. We specify their
hypotheses and prove the replacements that concern the law of the field.
In particular, a result for Gaussian interfaces is not used as a theorem
about the present field.

\subsection{Local insertion, signs, and crossings}

\begin{lemma}[Bounded local insertion]\label{g:insertion}
For every fixed \(L<\infty\) there is \(p_L>0\) with the following property.
Condition on any admissible exact oriented strands. In a patch of diameter
at most \(L\), at distance at most \(L\) from one of their tips or from
bounded Dirichlet boundary, prescribe a compatible sign pattern on at most
\(L\) additional vertices. The probability of that pattern is at least
\(p_L\). In particular every feasible prescribed extension of at most
\(L\) interface steps has a probability bounded below by a constant
depending only on \(L\). There are also absolute \(p_*,b_*>0\) such that
every nonfixed positive string site has conditional probability at least \(p_*\) to
have height greater than \(b_*\), without any restriction on its distance
from strand tips or Dirichlet boundary. Both assertions remain valid when
other positive string sites have additional upper bounds in \((0,d]\),
provided neither a prescribed new positive site nor the tested string site
is one of those upper-bounded sites.
\end{lemma}

\begin{proof}
Every string site has an oppositely sign-constrained nearest neighbor,
so Lipschitzness on their edge bounds both heights in absolute value by one.
This supplies the local anchor for the stated patches and also for a
singleton patch at any string site, including one far from the tips.
Lipschitzness therefore bounds heights in the patch by a deterministic
\(M=M(L)\). Enclose it in
a larger fixed box and use Corollary~\ref{f:conditioned-sparse} with a long-distance
deficit greater than \(2M+20\); the small linear-slack version gives this
by enlarging the fixed box. The good-shell event has probability at least
\(1/2\), uniformly in the strand conditioning and the lengths of the
additional upper intervals.

Fix such a shell and give the prescribed vertices, including the tested
positive string site if present, their requested values
\(\pm b\), where \(0<b<1/4\). Clip any feasible filling between the lower
and upper distance cones of these prescriptions. At each other
sign-constrained vertex the clipping interval contains zero, since that
vertex has graph distance at least one from the prescribed vertices.
Consequently clipping preserves its sign. Since the clipping interval
contains zero, it also preserves every old interval \([0,d]\), without
making \(b\) depend on \(d\), and preserves
fixed exterior values of absolute value at most \(1/2\). The shell remains
unchanged because of the long-distance deficit. This is the local
feasible-filling construction underlying Lemma~\ref{f:strict-filling}.

Let \(g\) be the resulting filling. For every original filling \(h\), the
map \(h\mapsto t h+(1-t)g\), with
\(0<t<b/(2(M+b))\), gives the required signs on the patch and makes
every prescribed positive height greater than \(b/2\). Take \(b_*=b/2\)
for the singleton assertion. Convexity preserves all the old constraints. If the number of free box
coordinates is \(N\), its Jacobian is \(t^N\); both \(N\) and \(t\) depend
only on the fixed patch and box. Integrating the conditional volume bound
over good shells proves the result. For an interface extension prescribe
the signs on the two sides of its finitely many edges.
\end{proof}

\begin{lemma}[Comparison behind controlled pins]\label{g:screened-comparison}
Consider a one-sided experiment from Theorem~\ref{c:comparison} at
scale $R\to\infty$, after an adapted exact-value transcript has been
recorded. Let $G$ be a recorded union of components of the common
unqueried graph. Every original lattice edge from $G$ to its complement
must end at a common fixed vertex; no free--free edge is deleted in
selecting $G$. There are no unrecorded individual walls in $G$.
Write $P$ for all fixed vertices incident to $G$, $z=h-\widetilde h$,
and consider either
\[
 u=(a-z)_+,\quad P_{\rm c}=\{p\in P:z(p)\ge a\},
 \qquad\hbox{or}\qquad
 u=(z-b)_+,\quad P_{\rm c}=\{p\in P:z(p)\le b\},
\]
with bounded threshold. Put $P_{\rm u}=P\setminus P_{\rm c}$ and define
\[
 \delta_G(x,P_{\rm u})=
 \inf\{\operatorname{diam}_\infty\gamma:
 \gamma=(x,\ldots,p)\text{ is a lattice path with }
 \gamma\setminus\{p\}\subset G,\ p\in P_{\rm u}\}.       \tag{SC1}
\]
Fixed vertices are terminal, not traversable, in (SC1), and an empty
infimum is infinity.

Let $Q_{\rm in}\Subset Q_{\rm out}$ be coordinate squares fixed before
the transcript is sampled, with separation comparable to $R$ and
$\operatorname{diam}_\infty Q_{\rm out}\le(c_0-c_1)R$ for a fixed
$c_1>0$. Suppose that a tested set
$T\subset G\cap Q_{\rm in}$ satisfies
$\delta_G(x,P_{\rm u})\ge c_0R$ for every $x\in T$. Assume the original
annealed sparse-contact, hole, and deterministic-test average hypotheses
of Theorem~\ref{c:comparison} in a fixed enlargement of $Q_{\rm out}$.
Use full ambient Lipschitz extensions of both marginals. In particular
the negative Sobolev bound is for the full, unmasked $z$ there, including
its zeroth-scale anchored average. Then the one-sided conclusions of that theorem hold on
$Q_{\rm in}$ for the components meeting $T$, with the tail extended
by zero behind their controlled fixed boundary. Its hole extension $U$
satisfies, for every $1<p<2$,
\[
 \E\|U\|_{W^{1,p}(R^{-1}Q_{\rm in})}^2\le C,\qquad
 \E R^{-2}\sum_{x\in T}|U(x)-u(x)|\longrightarrow0.     \tag{SC2}
\]
It is zero at retained vertices of $P_{\rm c}$ in this window. The
constants may depend on the square ratios and $c_0$. The same statement
holds on a recorded event of probability at least $p_*>0$, with
constants depending on $p_*$. No estimate is asserted at a test with
short access to $P_{\rm u}$. A trace application must separately verify
thickness and transport to retained controlled pins; the artificial rim
is not a zero-tail trace anchor.
\end{lemma}

\begin{proof}
Inside $Q_{\rm out}$ retain the union $C$ of components of
$G\cap Q_{\rm out}$ meeting $T$. A fixed pin incident to $C$ before
the outer rim is controlled. Otherwise a path from one of its tests
to that pin, contained in $Q_{\rm out}$ apart from one terminal lattice
step, would have diameter less than $c_0R$ for large $R$, contrary to
(SC1). Every edge from $C$ to a removed component inside the square
therefore passes through a controlled fixed pin. Edges at the outer
rim are not re-pinned or deleted from the conditional baths: that rim
is only an artificial boundary for the local estimate. It is a fixed
positive ambient distance from $Q_{\rm in}$.
The theorem is applied once to this union, not separately with a constant
for each of its possibly many components.

These are precisely the local one-sided hypotheses of
Theorem~\ref{c:comparison}. To see the localization in its proof,
a transmission cluster which brings any uncontrolled input to an
inner cutoff must reach the outer rim. Its diameter is then a fixed
multiple of $R$, so it is covered by the large-cluster estimate after
\eqref{c:caccioppoli}. A fixed forcing box away from the rim has no
uncontrolled incident pin; each of its free components reaches its
artificial collar or a genuine controlled pin. Thus
Lemma~\ref{c:force} and the allocation argument use the original
annealed good marks. No estimate for a newly selected quenched pin
vector is being asserted.

For clarity, the tail is zero-extended only across the controlled
boundary inside $Q_{\rm out}$, where every new edge has a zero pin
between its retained and removed sides. This extension is uniformly
lattice-Lipschitz there. Do not set it to zero across the artificial
outer rim. The hole construction and the test cutoffs are taken in
nested inner squares with a fixed buffer before that rim. A hole
reaching from those squares to the rim has diameter comparable to $R$
and is in the arbitrarily small exceptional event of
Lemma~\ref{f:hole}. Off that event, a positive copied value $U(x)=u(y)$
comes from a retained rim vertex $y$ in $C$. The full-window Lipschitz
bound on the unmasked $z$ gives
\[
 U(x)^2\le(a-z(x))U(x)+C(1+d_x)U(x)
\]
for the lower tail, and the reflected inequality for the upper tail.
A copied zero satisfies it automatically, even on a removed component.
Pairing with the deterministic ambient cutoff is exactly
\eqref{c:one-sided-pointwise}--\eqref{c:one-sided-absorb}; it never
uses a randomly masked original field. The energy, Sobolev, and
original--filled estimates now give (SC2). Retained controlled pins
still have value zero. Conditioning on an event of probability at
least $p_*$ costs only its reciprocal in these annealed and
exceptional-event estimates.
\end{proof}

\begin{lemma}[Positive trace on a pure sign side]\label{g:sign-trace}
In a local sector with only positive string boundary, the limiting bounded
harmonic conditional mean has Brownian boundary trace at least \(c>0\),
away from endpoints and incompatible sides. Here a testing neighborhood
has a fixed access-diameter margin, in the actual free graph after the
string heights are frozen, from every incident boundary pin not controlled
by the comparison below; fixed pins cannot be traversed. On a negative sector its
trace is at most \(-c\). On an exterior boundary arc alone the trace is
its prescribed value. If the pure sector also has positive exterior pins
at least $\lambda_*>0$, its lower bound is $c\wedge\lambda_*$; the
negative mixed sector is reflected. These statements are sequentially uniform at
diverging lattice scales.
\end{lemma}

\begin{proof}
The conditional law is log-concave and monotone in its boundary data and
in its one-sided constraints. Lemma~\ref{g:insertion} gives a uniform
positive chance that a tested positive string height exceeds some
fixed \(b>0\), even after other positive sites have been restricted to
\((0,d]\). A one-dimensional log-concave density on \([0,1]\) with this
tail bound assigns at most \(C d\) mass to \([0,d]\). One proof integrates
the chord lower bounds for its logarithm between a point in \([0,d]\)
and a point carrying the prescribed tail mass. Thus, for small fixed
\(0<d\le1/2\), the tested height exceeds \(d\) with conditional probability at
least \(1/2\). Monotonicity makes the conditioning that all earlier
indicators vanish the worst one. Successive uniform couplings show that
the positive string heights dominate \(d\) times independent fair
Bernoulli variables.

Write $A$ for the full original string-height vector, and freeze all
its coordinates. Retain the actual free component or union of components
on the tested side. The positive portion in the preceding domination
may be \emph{all} positive string pins incident to this retained graph:
the successive bound is uniform in their number, and the newly tested
site is never one of the earlier upper-bounded sites. Put
$c_*=d/2$ for a string-only pure sector, and put
$c_*=(d/2)\wedge\lambda_*$ in a mixed sector. Using the same fair
indicators, couple pins $B_x\in\{0,2c_*\}$ below all these positive
coordinates of $A$, and extend to the remaining coordinates of $A$
by disintegration. Conditional on $(A,B)$, compare the original Gibbs
kernel on the retained graph with the kernel pinned to $B$ at these
positive string pins and to $c_*$ at every remaining incident pin.
The geometry of the retained graph is fixed by the strand observation,
not by $A$ or $B$. The comparison mixture is invariant under
$h\mapsto2c_*-h$, so its mean is $c_*$. Its pins lie in
$[0,2c_*]\subset[-1/2,1/2]$ and have the uniform comparison estimates.

Rejoin the conditional kernels synchronously, keeping the conditional
baths on every other original component. Lemma~\ref{c:freezing} then
preserves the full original annealed marginal. Extend the comparison
field by $c_*$ off the retained graph, except at its $B$ pins; this is
a uniformly Lipschitz bounded-pin extension. Thus deterministic ambient
tests of both full fields have their required anchored and zero-mass
moments. For the lower tail $u=(\widetilde h-h)_+$, all selected string
pins are controlled because $A_x\ge B_x$. Every incident positive
exterior pin in the mixed sector is controlled because it is at least
$\lambda_*\ge c_*$. All remaining incident pins failing this inequality
are included in $P_{\rm u}$ of Lemma~\ref{g:screened-comparison}; none
is silently discarded.

The access hypothesis permits a fixed finite cover of a testing
neighborhood by inner squares whose outer diameters are smaller than
that margin. In each square take $T$ to include every free vertex under
the smooth tests in the screened neighborhood. Lemma~\ref{g:screened-comparison}
gives the filled lower tail with uniformly bounded Sobolev norm on this
cover, and (SC2) identifies its pairings with the original tail there.
Exhausting such test supports identifies the distributional response
only on the open screened side, which is all that is needed.

For the trace, choose a component $C_0$ of $G\cap Q_{\rm out}$ meeting
$T$, and a compact accessible pure subarc of the boundary of its
free-vertex hexagon union, contained in a smaller square compactly
inside $Q_{\rm in}$, of diameter at least a fixed multiple of $R$.
Trace that hexagon boundary along the subarc. The neighboring
fixed hexagons are incident pins; consecutive ones meet at a tiling
vertex and their triangular-lattice centers are adjacent. The local
access condition makes all these pins controlled. They therefore give
a connected controlled chain of diameter comparable to $R$. This chain
has the required local thickness, and its vertices can be moved past
the polylogarithmic holes to retained controlled pins. Apply the
localized proof of Proposition~\ref{t:varying-trace}, with $3/2<p<2$,
only to measures on a smaller compact part of this chain. The Sobolev
square is its buffer; the artificial rim is not used as an anchor.
It gives zero trace for every limit of the filled tail mean there.
Together with the interior identification above this proves the needed
tail inequality. No harmonic equation for the tail is used.
Jensen's inequality bounds the negative part of the limiting mean
difference by this nonnegative zero-trace limit. The two field means are
bounded, harmonic away from the observed sets, and their remote
contribution is bounded by a constant times the harmonic probability of
reaching an endpoint or an incompatible side. This probability tends to
zero along Brownian approaches to the pure portion, by Beurling.
The lower trace is therefore at least \(c_*\). All bounds used are
uniform for each deterministic admissible conditioning sequence, which
gives the claimed sequential uniformity. Reflect the field for the
negative statement.

The graph-access condition has a concrete meaning for the uses below.
After freezing all strand-adjacent heights, every primal edge crossed
by a dual strand has both endpoints fixed. A free path cannot cross that
strand without first ending at one of those pins. A nearby negative pin
on the other side may still be incident to the global component of an
open prefix, but it is not reachable from the tested side by crossing
the slit; any remaining route to it must go around an endpoint or reach an
incompatible boundary, and the asserted access margin is exactly what
excludes such a short route. In the endpoint collars this margin is
verified below. For the completed reference component with exterior
values $\pm1/2$, take all incident positive string pins in the coupling.
Every other incident pin on its positive component is $+1/2\ge c_*$,
and no negative pin is incident: the crossed-edge incidence also closes
the two marked ends, as detailed in Lemma~\ref{h:terminal-sign}.
Here $P_{\rm u}$ is empty and its access diameter is infinite. This
completed-component application, not an open-tip comparison, is the
strict positivity input for the reference height gap.
\end{proof}

\begin{lemma}[The complete-crosscut upper bound]\label{g:upper-half}
For a complete single crosscut, the positive-side conditional mean is
at most \(1/2\), and the negative-side conditional mean is at least
\(-1/2\). This assertion is not made at an open interface tip.
\end{lemma}

\begin{proof}
Condition on all heights on the negative side. At a positive string
vertex, a crossing edge gives an upper bound at most one. Delete the
crossing edges and replace these bounds by independent walls \([0,1]\).
This raises the conditional field by monotonicity. Raising the exterior
positive data to \(1/2\) raises it again. The resulting positive-component
law is invariant under \(h\mapsto1-h\), so its mean equals \(1/2\).
Average the comparison over the conditioned negative heights. The
negative statement follows by reflection. Disconnection of the two
components by the full crosscut justifies deletion of precisely these
crossing constraints.
\end{proof}

\begin{lemma}[High crossings]\label{g:high-crossing}
In a compact corridor of fixed shape and size \(n\), with an observed
component of comparable size and distance, the probability of a path
of heights at least \(K\) crossing a fixed positive fraction of that
corridor tends to zero as \(K\to\infty\), after \(n\to\infty\).
The same is true for the difference of two independent samples under
possibly different bounded exterior or strand conditionings.
In both statements the corridor has an enlarged neighborhood of
width \(cn\), for a fixed \(c>0\), containing no original fixed pins
or sign walls. In the two-field statement this neighborhood lies
in both domains with the same unconditioned lattice graph, and
each original law has the stated macroscopic anchor. The two
conditionings may differ outside this neighborhood.
\end{lemma}

\begin{proof}
We give the two-field proof. Explore every cluster of
\(h_1-h_2\geq K\) connected to a deterministic transverse seed.
Every queried vertex away from the seed has difference at least \(K-2\):
positive queries satisfy the threshold, and failed queries are adjacent
to a positive query. Freeze both samples on the common queried set,
then use the stationary synchronous joining of their conditional laws.
The exploration decisions used only recorded values; hence this
replacement preserves each annealed marginal. All statements below are
annealed over this transcript, not asserted for every possible
conditional pin value.
Use a middle testing subcorridor a fixed positive distance from
the transverse seed. The free-buffer hypothesis removes all
original pins and sign walls from its larger comparison window;
the newly queried common pins there satisfy the displayed
threshold inequality. Thus the one-sided comparison has no
uncontrolled original boundary in that window.

Away from the seed, set
\[
 u_K=\left(1-\frac{h_1-h_2}{K-2}\right)_+ .
\]
It has zero data on the revealed set. The local truncated comparison
estimate in Theorem~\ref{c:comparison}, applied after division by
\(K-2\), gives compact filled extensions with uniformly bounded local
Sobolev energy. Here the cutoff energy grows at most quadratically with
the truncation level: after division it is bounded by the squared cutoff
norm and the macroscopic second moments of \(h_1/(K-2)\) and
\(h_2/(K-2)\). Proposition~\ref{f:averages} controls these moments uniformly.
The same bounds hold after conditioning on an event of probability
bounded below, with its reciprocal as a constant.

Suppose a crossing had probability bounded below along \(K\to\infty\)
and \(n\to\infty\). On that event the revealed set contains a connected
continuum of fixed positive diameter away from the seed. Pass to a
subsequence jointly with this continuum and the filled extensions.
Macroscopic averaged field tightness gives
\((h_1-h_2)/(K-2)\to0\) as a distribution in this iterated limit.
Since \(u_K\geq1-(h_1-h_2)/(K-2)\), its limit is at least one almost
everywhere. The varying-set anchor argument of
Proposition~\ref{t:varying-trace} gives zero trace for that same limit.
Here one uses its samplewise inequalities, without a Poisson equation
for the truncated tail. Choose a probability measure on the crossing
continuum by coordinate projection, with ball masses bounded by $Cr$
in physical units. Move it along the revealed set past each hole to a
retained zero vertex. The maximum hole diameter is polylogarithmic
outside the negligible exceptional event, and the crossing component
cannot be contained in such a hole. Lemma~\ref{t:frostman} bounds the
integral of a radius-$r$ average of the filled tail by
$C(r^\beta+o(1))\|U_n\|_{W^{1,p}}$, for some $\beta>0$.
These inequalities hold for the random measure before expectation.
Joint compactness and strong local $L^p$ convergence permit the mesh
limit at fixed $r$; the bounded conditional Sobolev norm permits
$r\downarrow0$. This proves the zero trace. A \(W^{1,p}\) function which is
at least one almost everywhere cannot have zero trace on a continuum
carrying a fixed positive-diameter Frostman test with exponent
greater than \(2-p\). This is a contradiction.

For one field retain the entire original field on the deterministic
testing window. Given its transcript, compare its unvisited Gibbs
kernel with a symmetric field zero-pinned on the revealed set and
original exterior, and extend this comparison field by zero on its
fixed coordinates. Keep the original field, not its restriction to
a random unvisited component, in the full-window difference.
For each recorded set, deterministic ambient tests of the comparison
field have the bounds of Proposition~\ref{f:averages}; the revealed
crossing component provides a macroscopic anchor on the crossing
event. The original full field keeps its annealed bounds by
disintegration. Queried nonseed differences are at least \(K-1\),
and the same lower tail with \(K-1\) has zero input there. The
free buffer excludes all original walls from this comparison window.
The preceding argument therefore applies without a random mask
on the original field.
\end{proof}

\subsection{Barriers and exact strand comparison}

For a union \(K\) of observed strands, let \(D(K)\) be the domain obtained
by deleting closed lattice triangles meeting \(K\).  When an endpoint
lies on the boundary of an open set \(U\), retain its specified incident
access.  Define
\[
 \begin{aligned}
 d_U(A,B)=\inf\{\operatorname{diam}\alpha:\;&
   \alpha:[0,1]\longrightarrow\overline U\text{ is continuous},\\
   &\alpha((0,1))\subset U,\quad \alpha(0)\in A,\quad\alpha(1)\in B,\\
   &\text{with the specified endpoint accesses}\}.
 \end{aligned}
\]
For a barrier endpoint its access is the one selected by its initial
prefix.  Thus distinct slit banks or free sectors meeting only at a
cellular pinch are not identified.  The distinction between this internal
diameter distance and Euclidean distance is essential below.
All assertions below that a barrier pocket has a simple boundary arc, or
that a sign crosscut joins two pure seed incidences and separates the
pocket, are assertions on this access-resolved surface.  If the same fixed
pin occurs at several boundary incidences, those incidences are distinct
terminal boundary points but retain the one common fixed height.

A \((\theta,R)\)-barrier is a path \(\Upsilon\) of diameter in
\((\theta R,R]\), with an observed or exterior boundary component of
diameter at least \(\theta R\) within distance \(\theta^{-1}R\).
Its interior is in \(D(K)\). At each endpoint \(z\) on the boundary,
the first encountered circular crosscut at radius \(\theta R\) bounds
a pocket \(A_z\) with one simple boundary arc; the path meets every
smaller centered circle once inside this pocket. If a point of the
barrier is within \(\theta R/5\) of the observed boundary, it belongs
to one of these pockets and lies within \(\theta R/2\) of its endpoint.
These are the four geometric conditions of
\cite[Section 3.4]{SSdiscrete}. A signed barrier is compatible if the
boundary arc of each endpoint pocket has the indicated sign.
Oppositely signed barriers must have internal distance at least
\(2\theta R\). At bounded lattice scales no hexagon may meet both
prescribed signs.

\paragraph{Straight-attachment layouts.}
Put \(U=D(K)\) and \(d=\theta R\).  All clearances here are from the
full conditioned boundary \(\partial U\), including the exterior.
The diverging-scale barrier layouts used below have a fixed \(a>0\)
with the following additional properties.  At each boundary endpoint
\(z\), the initial part through radius \(d\) is a straight nearest
attachment
\[
 \gamma_z(r)=z+r e_z,\quad 0\le r\le d,\qquad
 \operatorname{dist}(\gamma_z(r),\partial U)=r.
\]
It is the radial segment in its Jordan pocket \(A_z\).  If
\(\Upsilon^{\rm f}\) is the barrier with only the initial portions
\(\gamma_z([0,3d/5))\) removed, then
\[
 \operatorname{dist}(\Upsilon^{\rm f},\partial U)\ge ad,\qquad
 \gamma_z(31d/32)\in\Upsilon^{\rm f}.                            \tag{UB1}
\]
Interior endpoints are part of \(\Upsilon^{\rm f}\).  A bounded
lattice error in these inequalities is harmless at diverging \(d\).

The proof constructs confined attachment domains and deterministic
crossing rectangles with fixed free interpolation neighborhoods.
Let \(C_s\) be the union of the full confined domains of sign \(s\),
with one lattice neighbor layer, and let \(P_s\) be the union where
the sign-\(s\) prescription is nonzero, including its interpolation
neighborhoods.  For a mixed layout require, for a fixed \(\kappa>0\),
\[
 d_{\rm f}(C_+\cup P_+,P_-)\ge\kappa R,\qquad
 d_{\rm f}(C_-\cup P_-,P_+)\ge\kappa R.                          \tag{UB2}
\]
Here \(d_{\rm f}\) is the infimum of the diameters of paths in the
actual free graph between the indicated sets; fixed or grounded
vertices are terminal and cannot be traversed.  The choices are made
from the conditioned geometry before sampling heights.  For a
single sign (UB2) is empty.  It does not separate \(C_+\) from \(C_-\).

\begin{proposition}[Uniform barriers for straight-attachment layouts]
\label{g:barriers}
Fix \(\theta>0,m<\infty,a>0\), and for mixed layouts fix \(\kappa>0\).
There is \(p>0\) such that, for every compatible collection of at most
\(m\) \((\theta,R)\)-barriers with the preceding layout, the conditional
probability that all continuations of the observed strands avoid all
the barriers is at least \(p\).
Pure exterior attachments are allowed when the prescribed heights are
at least \(\lambda_*>0\) on a positive attachment and at most
\(-\lambda_*\) on a negative attachment; \(p\) may then depend on
\(\lambda_*\).  Interior attachment constants do not depend on it.

For diverging-scale applications assume local thickness in the endpoint
testing windows: after cutting the observed system by a generous square
of side comparable to \(d\), each retained component meets its rim or
has diameter at least \(cd\), and a nearby attachment component has
diameter at least \(cd\).  At bounded lattice scales the compatible
finite sign fences are required not to prescribe both signs on one
hexagon; they are treated by insertion.
\end{proposition}

\begin{proof}
We give the diverging-scale proof.  All sets below are fixed by the
conditioned geometry.  Write \(c_{\rm lat}=O(1)\) for the lattice
boundary error.  Fix the constants in the cellular thickenings and
rectangle enlargements below, bounded by \(L\), then take
\[
 t=\eta R,\qquad \ell_c=t/M,
\]
where \(M\) is a sufficiently large fixed number.  Choose \(\eta\)
small and then \(R\) large so that
\[
 (L+C_0)\ell_c+c_{\rm lat}\le t/4,\qquad
 12t\le ad/256,\qquad L\ell_c\le d/128.                               \tag{UB3}
\]
Increasing fixed constants or decreasing \(a\) to at most \(1/2\)
does not change the statement.  Let \(Q\) be the union of the closed
coordinate cells of side \(\ell_c\) meeting the \(t\)-neighborhood of
all of \(\partial U\).  A non-\(Q\) cell is free and has distance
at least \(t\) from that boundary.  At a \(Q\)-facing incidence its
distance from a pin is at most \(2t+C_0\ell_c+c_{\rm lat}\).

\paragraph{The multi-face core and its tests.}
Fix a boundary endpoint and one sign side.  Cutting the Jordan pocket
\(A_z\) along the whole straight segment \(\gamma_z([0,d])\) gives
two Jordan sides.  Let \(P\) be the chosen one.  Its only genuine
conditioned boundary is one branch \(\beta\) of the pure simple
boundary arc, running from radius \(0\) to radius \(d\).
Let \(V\) be the component of
\[
              P\cap\{5d/8<|x-z|<15d/16\}
\]
adjacent to the cleared segment.  Its genuine boundary pieces \(B\)
lie on \(\beta\).  Its artificial boundary \(J\) consists of the
two circle cuts and the radial cut.  Indeed the boundary of a
component of this open intersection lies on those stated boundaries.

Here is a literal fine-cell convention.  In the original cut
triangular-lattice graph admit every original free vertex whose
hexagon center is in \(V\), with a fixed tie rule on a cut.  Let
\(E_V\) consist of the original edges between admitted vertices
whose embedded edges do not cross a circle or radial cut, and write
\(\mathcal V\) for this confined graph.  It need not be connected
or a disk.  Every contact whose embedded edge crosses an artificial
cut is a \(J_\triangle\) incidence, even if both endpoint centers
are admitted; such an edge is omitted only from \(E_V\), not from
the original ambient graph.  For an original edge from an admitted
vertex to a nonadmitted vertex, classify its first crossing of
\(\partial V\) as \(B_\triangle\) when it is a verified fixed
incidence on the pure branch, and as \(J_\triangle\) when it crosses
a circle or radial cut.  A cut crossing takes precedence at a corner.
The \(J_\triangle\) contacts are within \(c_{\rm lat}\) of \(J\).
A free--free edge cannot cross the conditioned branch in the original
cut graph: the strand-crossed primal edges have fixed endpoints.
Since no other interior vertex or edge is omitted, this classifies
every fine boundary contact, including fine rasterization holes.
There is no identification through a conditioned slit.  Eligibility,
target adjacency, and components below all use \(E_V\).  Exploration
uses \(E_V\) between free vertices and the verified \(B_\triangle\)
seed incidences for its initial queries.  The comparison graph always
retains the full original edge set.

Take \(b=ad/32\) and the parallel spine
\[
 \sigma(r)=z+\sqrt{r^2-b^2}\,e_z+b n,\qquad
        21d/32\le r\le29d/32,
\]
where \(n\) points into \(P\).  The nearest-segment clearance and
(UB3) put a fixed \(\ell_c\)-neighborhood of this connected spine in
\(V\setminus Q\), outside the closed \(10t\)-band about the radial
cut.  Call a free vertex eligible if it lies in \(\mathcal V\), its
coarse cell is non-\(Q\), and that whole cell is outside the \(10t\)
band.  Let \(S\) be the component of the induced eligible fine graph
containing the spine neighborhood.  The latter is connected at the
fine scale for large \(R\), so \(S\) is well-defined and connected.
It has a spine of diameter comparable to \(d\).  Its coarse covering
cells are non-\(Q\), stay outside the radial band, and number
\(O((d/\ell_c)^2)\).  All boundary faces of \(S\) will be used; no one
face is selected.

For each \(r\in[11d/16,13d/16]\), follow the circle of radius \(r\)
from \(\sigma(r)\) on the chosen side away from the radial bank
until its first hit \(q(r)\) of the closed set \(Q\).  This hit
exists.  The circle arc in the Jordan side reaches \(\beta\), hence
\(Q\), before the opposite radial bank.  Before the hit it is
non-\(Q\) and outside the excluded band.  The straight cut and its
clearance give this last assertion also near a possible return to
the other bank: a point of \(Q\) is reached before a path on this
side can enter that bank's cleared neighborhood.  With the cell
errors in (UB3), the pre-hit cells connect to the spine component.
At a grid corner an interposed \(Q\) cell is already a hit because
\(Q\) is closed.  Approximating this central arc by fine edges gives
an adjacent core cell whose middle vertices belong to the same \(S\).

Choose that cell by a fixed tie rule, move \(q(r)\) to its middle,
and smear there with a fixed smooth kernel of radius \(c\ell_c\).
Push normalized Lebesgue measure in \(r\) forward by this rule.
A Euclidean ball of radius \(\rho\) meets only radii in an interval
of length \(2\rho\).  The \(O(\ell_c)\) movement therefore gives
\(\mu(B(x,\rho))\le C(\rho+\ell_c)/d\).  Each cell receives at most
\(C\ell_c/d\) mass; the smearing has density at most \(C/(d\ell_c)\).
Consequently
\[
                  \mu(B(x,\rho))\le C\rho/d
                  \quad(0<\rho\le d).                            \tag{UB4}
\]
The hits may jump between arbitrarily many boundary faces.  For fixed
\(\eta\), the tests are a compact smooth family with support in \(S\),
because their cells lie in a bounded disk and number \(O((d/\ell_c)^2)\).

Every test has distance at least \(c(a)d\) from \(J\).  The circle
cuts have radial distance at least \(d/16-O(\ell_c)\).  If \(q(r)\) is
within \(d/32\) of the radial segment, its nearest segment point has
radius in the cleared central range.  Since \(q(r)\) has a pin
witness within \(2t+C_0\ell_c+c_{\rm lat}\), its distance from the segment
is at least \(ad-(2t+C_0\ell_c+c_{\rm lat})\).  Thus it is in fact at
distance at least \(\min\{1/32,a/2\}d\) from the radial cut, up to
the errors in (UB3).  The same holds for the smeared tests.

A shortest segment from the adjacent core cell to its first point of
\(\partial U\) has diameter \(O(t+\ell_c)+O(1)\).  Stop the entire route
at its first conditioned vertex before replacing it by a fine lattice
path.  It cannot reach \(J\) by the preceding Euclidean separation,
so its first pin is an incidence of \(B_\triangle\).  Its first exit
from \(S\) is one of the core boundary incidences used below.  This
gives short access to a pure pin without crossing a slit.

\paragraph{The original positive mean.}
Freeze all original strand-adjacent heights.  Let \(G_0\) be the
component of the full original free graph containing \(S\).  A short
test-to-pin path just constructed ends at a genuinely incident pure
pin.  Trace the fine hexagon boundary of \(G_0\) through that incidence,
stopping at the first \(J_\triangle\) contact or when its \(G_0\)-side
leaves \(V\).  Every adjacent outside
vertex before that exit is fixed: another free component cannot share
an original edge with \(G_0\).  The boundary classification makes all
these incidences pure \(B_\triangle\), until \(J_\triangle\) is reached.
Successive outside hexagons, and also successive \(G_0\)-side hexagons,
are equal or triangular-lattice adjacent.

This boundary has a segment of diameter \(c(a)d\) near the initial
incidence.  To see this, its \(G_0\)-side has an arm through \(S\)
to the macroscopic spine.  On the other side follow the incident
pure branch in a thin exterior collar of the cut Jordan side until
its first exit through an inner or outer circle.  That wire is only
geometric; its endpoint differs in radius from the central incidence
by at least \(d/16-O(t)\).  Both arms leave a fixed \(c(a)d\)-ball.
A small closed binary hexagon contour would trap one of them by the
Jordan theorem.  A contour ending at \(J_\triangle\) must also leave
that ball, by the Euclidean separation just proved.  Stop at its
first exit.  The resulting connected pure-pin chain has diameter
comparable to \(d\).  The initial short free path and the consecutive
\(G_0\)-side cells stay in the same slightly larger ball.  Hence all
these pins are incident to the same local component of \(G_0\) in
that ball which meets the test.

This verifies both the graph-access margin and the incident controlled
trace chain in Lemma~\ref{g:sign-trace}.  A short free path to any other
pin would have to leave \(V\) through \(J\); fixed pins are terminal.
For a mixed string/exterior branch the lemma uses fair pins
\(\{0,2c_*\}\), with \(c_*=(d_{\rm ht}/2)\wedge\lambda_*>0\);
for a string-only branch \(c_*=d_{\rm ht}/2\).  Here \(d_{\rm ht}\)
is its fixed insertion height.  Its trace inequality, Beurling's
remote-boundary bound, and the \(O(t)\) pin access give
\[
 \liminf_{R\to\infty}\E\langle\mu,h\rangle
       \ge c_*-o_{\eta/\theta}(1).                               \tag{UB5}
\]
The test in (UB4) is uniformly \(1/2\)-Frostman.  The stated local
thickness and anchor assumptions allow
Corollary~\ref{t:local-boundary-variance}, so
\[
 \lim_{\eta/\theta\downarrow0}\limsup_{R\to\infty}
       \operatorname{Var}\langle\mu,h\rangle=0.                  \tag{UB6}
\]
The inner limit is always at fixed \(\eta\).

\paragraph{Exhaustive attachment exploration.}
Put
\[
 C=\mathcal V\setminus S,\qquad
 \mathcal T=\{v\in C:\text{an edge of }E_V\text{ joins }v
                                      \text{ to }S\}.             \tag{UB7}
\]
Thus \(\mathcal T\) contains every free target neighbor on every
face, including artificial coarse-cell rims.  Incidences outside
\(\mathcal V\) are \(B_\triangle\) or \(J_\triangle\), not queryable
targets.  At a fine corner the target assignment in (UB7) wins
whenever that \(E_V\)-adjacency exists; any other cut-crossing contact
of the same vertex remains in \(J_\triangle\).
No original pin is adjacent to \(S\), by its non-\(Q\) buffer.

Explore in all of \(C\) every strict positive component attached
to every positive \(B_\triangle\) seed.  Query a vertex only from
an already discovered positive neighbor and continue until all such
neighbors have been queried, unless a positive vertex in \(\mathcal T\)
is found.  That is the event \({\rm Att}\).  Denote the reached
positive vertices by \(\mathcal P_+\), and the failed queries by
\(\mathcal F\).  Each failed query has height in \([-1,0]\).
The core \(S\) is never queried.  On \(E={\rm Att}^{\,c}\),
\[
 \mathcal P_+\cap\mathcal T=\varnothing,\qquad
 U_0=\text{the component of }\mathcal V\setminus\mathcal P_+
                                      \text{ containing }S        \tag{UB8}
\]
contains every target vertex.  Failed queries are not removed when
forming \(U_0\); some may lie in other complementary components.
Thus this is not the complement of the full query transcript.

Color \(U_0\) one color and all other fine cells of \(\mathcal V\)
the hull color.  Adjoin to the hull color a thin exterior wire along
every component of \(B\), joined to all its \(B_\triangle\) contacts
and continued along its pure branch up to \(J\).  A component in the
open annulus ends on a circle cut, even when both ends lie on the same
circle; tangencies and overlaps with a cut are assigned to \(J\).
The wires use fixed local detours at fine corners and are geometric
only: they add no Gibbs values or graph identifications.  In particular
a contour cannot terminate at a \(B_\triangle\) boundary corner.
For a short reversed pin-to-test path, take its last
hull-to-\(U_0\) edge \(vu\).  If \(v\in\mathcal V\), then
\(v\in\mathcal P_+\); otherwise it would be an adjacent vertex of
the same component of \(\mathcal V\setminus\mathcal P_+\).
Also \(u\notin S\), because a positive neighbor of \(S\) would be
a successful target.  Exhaustive exploration therefore queried \(u\)
from \(v\), and \(u\in\mathcal F\).  The same conclusion holds for
a direct edge from an actual \(B_\triangle\) seed.  The short path
does not meet \(J_\triangle\).  We have found a failed query within
\(O(t)\) of each test.

The binary fine-hexagon contour through this edge supplies a nearby
failed chain of diameter \(t\).  At an interior honeycomb vertex the
binary interface has degree two, regardless of the number of faces
or holes.  Along an \(E_V\) hull--\(U_0\) edge the hull neighbor
is reached positive by the argument above; a filled complementary
component cannot edge-border \(U_0\).  The \(U_0\)-side vertex is
therefore failed.  The same is true at a genuine seed edge; an edge
of \(J_\triangle\) is the only uncontrolled possibility.

There are two escaping geometric arms at the initial edge.  The
\(U_0\)-side connects in \(U_0\) through \(S\) to a spine point at
distance \(c d\).  The hull-side connects through its actual positive
cluster to one of its seeds.  Follow that seed's own pure branch in
the exterior collar of the cut side until its first exit from the
radial annulus.  A returning branch piece can have both ends on the
same circle, which suffices.  The branch runs from radius \(0\) to
radius \(d\), so the exit is on an inner or outer circle, not on
the radial cut inside this band.  Its radius differs from that of
the hit by at least \(d/16-O(t)\).  This is a Euclidean diameter
bound for the wire, not a free-graph access assertion along fixed pins.
Both arms can be placed strictly on their respective color sides.
If the contour closed in \(B(u,2t)\), the arm on its bounded side
could not escape, by Jordan separation.  It cannot encounter
\(J_\triangle\) in that ball, by the Euclidean margins above.
It therefore exits the ball.  Consecutive \(U_0\)-side cells are
equal or triangular-lattice adjacent, giving a connected failed
chain of diameter at least \(t\) within \(O(t)\) of the test.

\paragraph{Comparison on attachment failure.}
Include all original string heights and all queries in the exact
transcript.  Let \(G\) be the union of the whole components of the
actual unqueried graph meeting the tested vertices.  Compare the
original conditional kernel on \(G\) with the zero-pinned uniform
kernel \(h^0\) on \(G\), using a stationary synchronous joining and
retaining every other original conditional bath.
Lemma~\ref{c:freezing} preserves the full original marginal.
The comparison kernel has conditional mean zero.  Extend \(h^0\)
by zero off \(G\).  It is a full Lipschitz zero-pinned field, and the
nearby original attachment component is a connected zero-pin anchor
for this extension.  All negative-Sobolev estimates use the full
\(z=h-h^0\), not a masked original field.

For \(u=(h-h^0)_+\), all incident pins of nonpositive original height
are controlled; every positive original or queried pin is counted as
uncontrolled.  No unqueried path from a test can reach such a pin
when all its preterminal free--free steps use \(E_V\) and its terminal
incidence is not in \(J_\triangle\).  If the terminal pin is in
\(\mathcal P_+\), its predecessor in \(C\) was queried by exhaustive
closure, and a predecessor in \(S\) would make it a successful target.
For a \(B_\triangle\) pin, a predecessor in \(C\) was likewise queried,
and a predecessor in \(S\) is excluded by the buffer.  All positive
original pins incident before a \(J_\triangle\) contact are among
those seeds, by the fine boundary classification.  Hence every
ambient unqueried path to an uncontrolled pin must first cross
\(J_\triangle\), either by leaving the admitted vertices or by using
an excluded cut-crossing edge with admitted endpoints.  Its diameter
is at least \(c(a)d\) for every tested vertex.  This proves (SC1)
without restricting the comparison graph to \(\mathcal V\) and
without deleting any free--free edge.

Cover the tests and their \(O(t)\) neighborhoods by a fixed finite
family of preassigned inner squares whose outer diameters are smaller
than this access margin and whose inner-to-rim gaps are fixed multiples
of \(d\).  Apply Lemma~\ref{g:screened-comparison} once to the union
of local components meeting the tests in each square.  On a recorded
event with \(\Prob(E)\ge p_*>0\), for \(3/2<p<2\) its filled tails
satisfy
\[
             \E[\|U\|_{W^{1,p}}^2\mid E]\le C_{p_*}.            \tag{UB9}
\]
The constants are independent of \(\eta\).  The tail is zero-extended
behind controlled fixed boundaries inside the buffered windows, never
across the artificial outer rim, and copied hole values are paired
with the full \(z\).  The event restriction costs its reciprocal.

Partition \(\mu\) into squares of side \(t\).  Each has mass at most
\(Ct/d\).  Move that mass to one of the nearby failed chains of
diameter \(t\), using the coordinate-projection measure in
Lemma~\ref{t:transport}.  Each chain stays within \(O(t)\) of its
assigned square.  The lemma's packing argument gives a common
Frostman bound and displacement \(O(t)\).  Project to fine vertices
and move past the polylogarithmic holes along the chain.
At an inner-window vertex of a failed chain, either the vertex is
incident to the retained local union, in which case its nonpositive
height makes it a controlled fixed pin, or it lies outside that union,
where the pre-filling tail was zero-extended behind the controlled
boundary.  Thus these are trace zeros before hole filling; copied
hole values are handled by the preceding hole transport.  The whole
chain is not asserted to lie in \(P_{\rm c}\).
Choose the overlapping square cover so these movements stay inside
its buffered windows.  The samplewise trace and anchor estimates
of Lemmas~\ref{t:frostman}--\ref{t:anchors}, followed by (UB9), give
some \(\beta>0\) such that
\[
 \limsup_{R\to\infty}\E[|\langle\mu,U\rangle|^2\mid E]
              \le C_{p_*}(t/d)^{2\beta}.                        \tag{UB10}
\]
The original--filled estimate transfers the first moment at fixed
\(\eta\).  Since \(h\le h^0+u\) and \(h^0\) has mean zero conditional
on the transcript,
\[
 \limsup_{R\to\infty}\E[\langle\mu,h\rangle\mid E]
              \le C_{p_*}(t/d)^\beta.                           \tag{UB11}
\]
For a random variable of variance \(v\),
\[
 |\E[Y\mid E]-\E Y|\le\sqrt{v/\Prob(E)}.
\]
Equations (UB5)--(UB6) and (UB11), first taking \(\eta/\theta\) small
and then \(R\) large, exclude \(\Prob(E)\ge p_*\).
Thus the failure probability of each full multi-face attachment can
be made arbitrarily small.  Reflect for the negative side.

\paragraph{Protectors, the tube, and the common gate.}
We recall the elementary cellular crossing rule used throughout.
Place a longitudinal rectangle along each edge of a cell path and
a square at each vertex.  Require both coordinate crossings in the
square; a longitudinal edge crossing traverses that square and meets
its transverse crossing.  Concatenation gives a connected lattice
path.  Along a simple closed cell curve in a regular annular strip,
the concatenated walk has degree one under retraction to the curve.
Its decomposition into simple lattice cycles contains a winding-one
cycle.  These are literal common vertices, not mere overlaps of supports.

For each free part in (UB1), let \(\mathcal I_{\rm f}\) be its
coarse covering cells and set
\[
 N=\bigcup_{q\in\mathcal I_{\rm f}}
                 \bigl(q+[-5\ell_c/4,5\ell_c/4]^2\bigr),
\]
where \(q\) denotes the cell center.  The nonintegral half-width
removes corner-only contacts, with a fixed local tie convention.
There are \(O(\eta^{-2})\) cells per barrier, and the boundary
components are polygonal Jordan curves with feature size \(c\ell_c\)
and \(O(|\mathcal I_{\rm f}|)\) edges.  Free-part points have
coordinate distance at least \(3\ell_c/4\) from \(\partial N\); all
observed tips are outside \(N\) by (UB1).  On the \(N\)-side of
every boundary component require a winding-one sign cycle in a
regular strip.  Its fixed enlargement is free.  We do not fill
the holes of this tube: a complementary hole containing a tip has
its own cycle, separating that component from the buffered free part.

For a side core \(S\), let \(\mathcal I_S\) be all coarse cells
containing its vertices.  This covering is connected up to corner
contacts because \(S\) is connected in the fine graph.  Every target
neighbor is within one lattice step of it.  Take three nested regular
square-cell neighborhoods \(N_0\Subset N_1\Subset N_2\) of
\(\mathcal I_S\), with gaps \(c\ell_c\), total thickening at most \(L\ell_c\),
and with every target neighbor in the interior of \(N_0\).
Let \(H_j\) fill the bounded planar holes of \(N_j\).
Their boundaries have \(O((d/\ell_c)^2)\) edges and remain within \(L\ell_c\)
of the core covering.  Since that covering is non-\(Q\) and outside
the radial band, (UB3) puts these boundaries and their enlargements
inside the Jordan side \(P\), free of every pin and of the radial cut.
They have radius at most \(15d/16+L\ell_c\) and at least \(5d/8-L\ell_c\).
Filling preserves containment in \(P\) and in the outer disk: the
complement of each has a path to infinity.  Pure seeds escape along
their boundary branch outside the neighborhoods.  Thus the filled
holes contain no conditioned tips.

Require a sign cycle in a regular strip about \(\partial H_1\)
inside \(H_2\setminus H_0\).  It winds once around the whole \(H_0\).
Every successful attachment starts at a pure seed outside \(H_2\)
and ends at some target neighbor inside \(H_0\), so it meets this
cycle, whichever core face it used.

Move the common gate to \(g_z=\gamma_z(31d/32)\).  It is outside
the whole protector hull, since \(31d/32-(15d/16+L\ell_c)\ge3d/128\).
Join a spine cell inside \(H_0\) to that station through the cleared
parallel radial corridor, using non-\(Q\) cell rectangles on each
side.  Their count is \(O((d/\ell_c)^2)\); their enlargements remain
free, in the pocket, and below radius \(d+L\ell_c\).

For the shared crossbar use the component of the circle of radius
\(31d/32\) in \(A_z\setminus Q\) which contains \(g_z\).
Its two ends are adjacent to \(Q\), so their boundary distance is
at most \(2t+C_0\ell_c+c_{\rm lat}\).  The unrestricted tube \(N\) lies
within \(L\ell_c\) of the globally clear free part, and hence has boundary
distance at least \(ad-L\ell_c\).  By (UB3) both crossbar ends are outside
\(N\).  Choose a square \(G_z\) of side \(\ell_c/4\) about \(g_z\);
its closure is inside the buffered interior of \(N\) and outside
\(H_2\).  Its fixed enlargement is free.  Split the crossbar at
\(G_z\), terminate both radial connectors in it, and require both
coordinate crossings there.  The longitudinal bridge from each of
these paths traverses \(G_z\), so the corresponding random paths
meet its transverse crossing inside \(N\setminus H_2\).
The crossbar may meet several raw tube banks.
Perturb its cell path transversely to the bank edges and, at every
encountered bank, require a transverse crossbar crossing and a
longitudinal bank crossing in the same square.  They intersect
literally.  The common crossings in \(G_z\) join both radial
connectors to both crossbar halves.  Unencountered hole cycles may remain separate;
they still screen their own complementary components.

Each forced connector starts inside \(H_0\) and reaches \(G_z\)
outside \(H_2\), so it also intersects the winding-one protector.
Thus each full attachment joins the common gate and the encountered
tube banks.  Every endpoint module, with its fixed enlargements,
lies in \(B(z,d+L\ell_c)\).  The total number of crossing requirements
is \(O(m\eta^{-2})\), independent of microscopic boundary length.

\paragraph{Forcing and screening.}
Choose \(\eta\) so the attachment failures sum to less than \(1/10\)
for large \(R\).  In a positive rectangle planar duality gives a
crossing of heights \(>-H\) unless there is a transverse crossing
of heights \(\le-H\).  Lemma~\ref{g:high-crossing} excludes the
latter as \(H\) increases; reflect for negative rectangles.
Every rectangle has a fixed free enlargement of size comparable to
\(\ell_c\).  A subpiece of the nearby macroscopic anchor supplies its
anchor, with constants depending only on the fixed parameters.
Choose fixed \(H\) so all threshold failures sum to less than \(1/10\).
All threshold crossings and attachments then occur together with
probability at least \(4/5\).

Choose deterministic bumps equal to \(H+1\) on the positive crossing
rectangles and squares, equal to \(-(H+1)\) on the negative ones,
and zero outside their free interpolation neighborhoods and at all
grounded sites.  Maxima of same-sign bumps are allowed.  Their magnitude
is bounded, their edge differences are at most \(C/R\), and their
support has \(O(mR^2)\) vertices.  Condition (UB2) prevents conflicting
nonzero bumps.  Apply Proposition~\ref{g:shifts} once to the joint event.
On its retained subevent the corrected shift differs from the
prescription by \(o(1)\), turning all threshold crossings into strict
sign crossings.

The corrected shift also preserves all successful attachments.  A
near-contact component meeting \(C_+\) is either grounded, with zero
shift, or is a free component of diameter \(o(R)\).  In the latter
case (UB2) prevents it from meeting \(P_-\), so every prescription
averaged there is nonnegative.  The negative case is reflected.
A component meeting both collars but no prescription has zero shift.
All observed and seed sites are grounded.

For every resulting configuration, a continuation starting in a
complementary component of \(N\) cannot reach the buffered free part:
it would cross that component's strict same-sign cycle.  A dual zero
interface cannot cross a primal edge whose endpoints have the same
strict sign.  Every observed tip is outside \(N\).

At an endpoint the two attachments start on the two pure boundary
branches and join the common crossbar.  Select a simple path in this
joined sign graph from the last visit to one branch to the first visit
to the other.  Its interior lies in \(A_z\).  The attachment paths
stay in their sides \(P\); the protector paths have radii at least
\(5d/8-L\ell_c>3d/5\), and the connectors and crossbar are farther out.
The path is therefore disjoint from the initial prefix
\(\gamma_z([0,3d/5])\).  It is a crosscut of the Jordan pocket.
The side adjacent to the pure boundary arc containing \(z\) contains
that prefix and no observed interface tip.  A continuation reaching
the prefix would have to cross the strict sign crosscut.  Together
with the tube screening and (UB1), this covers every point of every
barrier.  No random curve is identified with a deterministic cell rim.

The feasible-shift proposition gives a uniform positive probability
for this image event.  All fixed constants, \(\eta\), and \(H\) precede
the limit \(R\to\infty\).  At bounded scales the compatible finite
sign fences lie in bounded patches at bounded distance from the anchor;
Lemma~\ref{g:insertion} supplies a positive lower bound there.
Take the smaller of the two bounds.

For the local-thickness hypothesis in the applications, external
arms and adapted continuations are attached to the original exterior.
An internal record is a connected incident strand through the sampling
center, of diameter comparable to its record radius at positive quality.
An endpoint testing square can be chosen not to contain that center.
If a smaller internal record meets the square, its connection to the
center crosses the square rim; otherwise its diameter is already
comparable to the square scale.  In the extra-excursion construction
the endpoint squares lie near radius \(\rho\), have size at most
\(c\rho\), and miss the center even when the inner record radius is
smaller.  The strand-swap comparison annulus is free of the inner
observations.
\end{proof}

\paragraph{Circular-clearance layouts.}
We use the one-sided circular construction in
\cite[proof of (3.48)]{SSdiscrete} with the following choices.
In its wide case stop the circular middle at boundary distance
\(b=\delta\rho/10\) and attach its ends to the two facing pure
pieces by globally shortest segments of common length
\(\ell_{\rm att}=b\).
In its \(k\)-th narrow case put \(b=\delta^k\rho\); the output
consists of two globally shortest segments from their common point
to the facing pieces.  If their common length is \(\ell_{\rm att}\),
the internal separation \(\delta b\) of those pieces and the route
through that point give \(\ell_{\rm att}\ge\delta b/2\).
In either case choose the endpoint radius
\(d=c_A\ell_{\rm att}\) with one fixed \(0<c_A\le1/(100A)\),
where the fixed \(A\ge1\) absorbs
the pocket and rectangle constants.  Along a shortest segment from
an endpoint \(z\), the distance of its radius-\(r\) point from the
full boundary is exactly \(r\), by the triangle inequality at its
other end.  After deleting the portions \(r<3d/5\), the remaining
segments have clearance at least \(3d/5\); the wide middle has
clearance \(b\), and the narrow output has no other part.
Successively exposed strands lie beyond the separated preceding
annulus.  Thus a fixed further shrink, followed by the bounded
lattice error, gives (UB1) with \(a=1/2\) at every diverging output
scale.  These are single-sign layouts, so (UB2) is empty.  Bounded
outputs use insertion.  This qualification applies separately to
each pure access used below.

\begin{lemma}[Exact-strand comparison]\label{g:strand-comparison}
Fix concentric inner and outer disks with radii in a fixed ratio
strictly greater than one. Consider two base conditioning systems
whose domains contain the closed outer disk and have the same
unconditioned lattice graph there. All vertices constrained by the
original strand observations, additional walls, or Dirichlet data
must be outside the outer disk in both systems. Thus the entire outer
disk, not just its annulus, is free in the two base laws. Each system
has anchored observed or Dirichlet pieces of comparable diameter
and distance outside this disk. Let \(I\) specify a bounded collection
of oriented strands all of whose adjacent sign-constrained vertices
lie in the inner disk. Then
\[
 C^{-1}\Prob_2(I)\leq\Prob_1(I)\leq C\Prob_2(I),
\]
with a constant depending only on the fixed ratios and number of
strands. The graphs need agree only inside the outer disk.
The assertion does not compare laws with different pre-existing
inner observations.
\end{lemma}

\begin{proof}
Take two independent samples from these base laws, conditioning the
first additionally on \(I\). The only observations inside the outer
disk are now this additional event in the first replica, and their
support lies inside the inner disk. Hence the annulus remains free
in both replicas. Each corridor used there has an anchor at comparable
distance under each law; an inner strand from \(I\) is also an
allowed anchor for the first replica.
Choose the two circuit supports in a closed middle subannulus
and cover them by finitely many corridors whose fixed
enlargements remain in that annulus. They satisfy the
free-buffer hypothesis of Lemma~\ref{g:high-crossing} in
both replicas. Lemmas~\ref{g:high-crossing}
and \ref{g:shifts} therefore give a
uniformly positive probability of two nested circuits in the free
annulus on which \(h_1-h_2\) has opposite strict signs.
We use the cluster-swapping construction of Sheffield
\cite[Chapter~8]{SheffieldSurfaces}, in the hard-constraint form
described below; see also \cite[Section~2]{CAP}.
For each edge declare it a dependency edge when swapping the
replicas at just one endpoint violates an edge constraint in one
of the replicas. This graph depends only on the unordered height
pairs. If the order of the replicas reverses across an edge, either
partial swap is feasible: each mixed endpoint difference lies
between differences of the original feasible samples. Hence a
dependency path cannot change the sign of \(h_1-h_2\).
No dependency component meeting the inner disk reaches outside
the two opposite-order circuits.

Swap both replicas at every vertex in a dependency component
meeting the inner disk. Every swapped component stays inside the
outer disk. The base laws impose no vertex observations there;
all their original observations therefore remain at untouched
vertices. A graph edge across a swapped-component boundary is
not a dependency edge, so the partial swap on that edge is feasible.
This proves preservation of both base polytopes. All vertices
tested by \(I\) are swapped, which transfers \(I\) from the first
replica to the second; the first replica is not required to retain
this additional event. The operation is an involution because the dependency
graph and selected components are invariant under swaps.
On each component-selection region it is a coordinate permutation,
so it preserves Lebesgue measure. If \(V_i\) is the volume of the
polytope with exterior conditioning \(i\), the uniform positive
fraction of swappable pairs gives
\[
 c\,V_1(I)V_2\leq V_1V_2(I).
\]
Reverse the experiments for the other inequality.
\end{proof}

\subsection{The harmonic-landing records}

Use rounded disks \(\mathfrak B_R\) equal to the union of the closed
Voronoi hexagons which meet the Euclidean disk \(B(0,R)\).  External
records stop at their first entrance to this union and reversed inner
records at their first exit.  A hexagon has circumradius \(1/\sqrt3\)
and diameter \(2/\sqrt3\).
Fix the origin and an incident dual edge \(e_\sigma^*\). Let \(S\)
be independent simple random walk from \(v\), let \(\tau_0\) be
its first visit to zero, and reverse its portion up to \(\tau_0\)
to obtain \(\breve S\). Let \(\tau\) be the first visit to a
vertex adjacent to the completed chord or to exterior boundary.
Define
\[
 Z^\sigma=\{S_\tau=0,\ e_\sigma^*\subset\gamma\}.
\]
Let \(\operatorname{ext}_R\gamma\) consist of the two initial
pieces grown from the marked endpoints until first entering
\(\mathfrak B_R\); it is the whole chord if the disk is missed.
Let \(\beta_R^\sigma\) be the part in \(\mathfrak B_R\) of the
zero interface through \(e_\sigma^*\), connected to that edge,
and put it equal to the empty path if there is no such interface.
Let \(\operatorname{ext}_R S\) end at the first entrance to
\(\mathfrak B_R\), and let \(\operatorname{int}_R\breve S\)
end at its first exit. Set
\[
 \Phi_R=(D,\partial_+,h_\partial,v,
              \operatorname{ext}_R\gamma,\operatorname{ext}_R S),
 \qquad
 \Theta_R=(\beta_R^\sigma,\operatorname{int}_R\breve S).
\]
The geometric path in \(\Theta_R\) is initially unoriented.
Conditioning on its orientation means specifying the adjacent signs.
Every subsequent field-law application below conditions on these
exact signs, then averages over the orientations.

Let \(x,y\) be the two external strand tips and \(q\) the external
walk tip at radius \(R\). When there has been no forbidden contact,
define
\[
 Q(\Phi_R)=\frac1{100}\wedge
 \frac{
 d(x,\operatorname{ext}_R S)\wedge
 d(y,\operatorname{ext}_R S)\wedge
 d(q,\operatorname{ext}_R\gamma)\wedge |x-y|}{R}.
\]
Put \(Q=0\) if the disk is missed by the chord or the revealed walk
hits a revealed strand. Define \(Q(\Theta_R)\) by the same formula
using the two inner strand tips and reversed-walk tip; put it equal
to zero for a closed inner interface or a forbidden walk contact.

Figure~\ref{fig:harmonic-records} depicts the two records and the
unexposed annulus between them.

\begin{figure}[htbp]
\centering
\begin{tikzpicture}[scale=.8,>=Stealth,font=\small]
  \fill[gray!10] (0,0) circle (3);
  \fill[white] (0,0) circle (1.15);
  \draw[dashed,gray] (0,0) circle (3);
  \draw[dashed,gray] (0,0) circle (1.15);
  \draw[very thick,blue!65!black]
    (.15,3.65) .. controls (-.1,3.4) and (.35,3.2) .. (.2,3);
  \draw[very thick,blue!65!black]
    (.1,-3.65) .. controls (.4,-3.4) and (-.15,-3.2) .. (.05,-3);
  \draw[very thick,blue!65!black]
    (.03,-1.15) .. controls (-.1,-.65) and (.2,-.4) .. (.15,.05)
    .. controls (.1,.5) and (.35,.8) .. (.16,1.14);
  \draw[very thick,red!65!black,->]
    (-3.9,.3) .. controls (-3.6,.1) and (-3.3,.3) .. (-3,.1);
  \draw[very thick,red!65!black,->]
    (0,0) .. controls (-.4,-.15) and (-.7,.18) .. (-1.14,.12);
  \fill (0,0) circle (.045);
  \node[below left] at (0,0) {\(0\)};
  \node[above left] at (-3.7,.32) {\(v\)};
  \node[anchor=west] at (1.35,.95) {\(\Theta_r\): inner record};
  \draw[gray] (1.3,.9)--(.65,.6);
  \node[anchor=west] at (2.4,2.65) {\(\Phi_R\): outer record};
  \draw[gray] (2.35,2.65)--(.35,3.2);
  \node at (1.8,-1.65) {buffer};
  \node[anchor=east] at (-.95,-1.8) {sampled side};
  \draw[gray,->] (-1.05,-1.55)--(-.45,-.4);
  \node[above right] at (.77,.77) {\(r\)};
  \node[above right] at (2.1,2.1) {\(R\)};
  \draw[very thick,blue!65!black] (-2,-4.15)--(-1.45,-4.15);
  \node[anchor=west] at (-1.3,-4.15) {observed interface};
  \draw[very thick,red!65!black] (-2,-4.55)--(-1.45,-4.55);
  \node[anchor=west] at (-1.3,-4.55) {observed walk pieces};
\end{tikzpicture}
\caption{The two records leave an annulus available for separation
and hookup. The inner walk arrow indicates reversal from the harmonic
landing. The shaded annulus is unexposed; no interface connection there
has yet been required.}
\label{fig:harmonic-records}
\end{figure}

\begin{lemma}[Hookup probability]\label{g:hookup}
For \(9R/8<R'<5R\), with the observer and exterior boundary outside
\(\mathfrak B_{6R}\), there is a universal upper bound for all
\(R\geq R_0\), where \(R_0>1\) is fixed:
\[
 \Prob(Z^\sigma\mid\Phi_{R'},\Theta_R)\leq C/\log R.
 \tag{G1}
\]
For every \(a>0\) there is \(c(a)>0\) such that the same probability
is at least \(c(a)/\log R\) when
\(Q(\Phi_{R'})\wedge Q(\Theta_R)\geq a\) and \(R\geq R_0(a)\).
\end{lemma}

\begin{proof}
The field part of \(\Phi_{R'}\) constrains only sites adjacent to
the external arms, hence outside radius \(R'-O(1)\).
The sites specifying \(\beta_R^\sigma\) lie inside radius \(R+O(1)\).
Choose comparison disks between these two radii, with
outer radius \((R+R')/2\) and inner radius \(R+(R'-R)/4\).
For sufficiently large \(R\), they have a free intervening annulus
of fixed relative width. The original base observations are all
outside the outer disk, and the external arms give the required
anchors. Thus Lemma~\ref{g:strand-comparison} applies to the two
orientations of the inner strand, comparing the exterior data with
their global sign reversal. Together with sign symmetry it gives a uniformly
positive probability that the two observed unoriented configurations
have compatible signs on the side containing the walk. In that case join corresponding interface tips and the two walk tips by
three polar centerlines.  We may assume \(a\le1/100\), since otherwise
the lower-bound hypothesis is empty.  After rounding their same-radius
angular gaps are at least \(a/2\).  Put
\(s=aR/(800\pi)\) and \(w=s/100\), reserve straight radial stems of
length \(20w\) at both ends, and interpolate the angles linearly between
the stems.  Since \(R'-R\ge R/8\) and \(40w\le R/(200000\pi)\),
the angular slopes are at most \(25\pi/R\).  At a common radius the
distance between centerlines is at least \(.49aR\).  If two radii differ
by at least \(s\), that gives the required separation directly; otherwise
moving along one centerline to the common radius costs at most
\((1+125\pi)s<.158aR\).  Thus the mutual separation is at least \(s\).
The observed external arms, the two free interface centerlines, and the
inner strand complete to one prospective simple oriented chord
\(\Gamma_0=K\cup J\), where \(K\) is the observed field system and
\(J\) its free completion.  The compatible orientation gives its two
sides \(\Omega_+\) and \(\Omega_-\).

Offset a lane by \(w\) on each side of each interface centerline.
At a tip the two offset lanes
meet its radius-\(4w\) circle at transverse offsets \(\pm w\) and
radial depth \(\sqrt{15}w\), up to bounded lattice error.  Follow from
the chosen lane the circular arc away from the radial stem.  The long arc
between the two lane points stays at distance at least \(w-O(1)\) from
\(J\); the other connector is farther by the centerline separation.
Its endpoints are on opposite ideal sides, so this arc must meet the old
boundary arc at that end.  At its start the full-boundary clearance is
greater than \(3w\), since the old record is on the old side of its
stopping circle and the other record is remote.  Stop at the first point
of full-boundary distance \(w/10\) and append a globally shortest segment
to \(b\in\partial D(K)\).  This segment remains in the selected side;
its endpoint is on the appropriate old pure bank, even if that bank
shadows the tip.  Moreover
\[
 \begin{aligned}
 d(b,J)&\ge .89w-O(1),\\
 d(b,\hbox{all observed tips})&\ge3.8w-O(1),\\
 d(b+tu,\partial D(K))&=t\quad(0\le t\le w/10).
 \end{aligned}
\]

Take \(d=w/(200A)\), with one fixed sufficiently large \(A\).  To check
the endpoint pocket, in each ideal side join free collars of the two
connectors \(J_1,J_2\) by the same-side arc on the middle-radius circle.
This arc is free and has macroscopic clearance, since the annulus has no
old observations.  Delete the triangles meeting \(K\) and use the
access-resolved component containing this connected reference set.  The
deleted triangles attach to its old boundary arcs, so its resolved
completion is a disk.  The radius-\(d\) crosscut met by the nearest
prefix misses the reference set.  The retained cap arc has clearance
\(w/10>d\), and a short transverse segment at its lane end reaches the
local radial stem with clearance greater than \(d\).  Thus the tail and
both connector collars are on the reference side of the crosscut.  The
other side is bounded by that crosscut and one pure old boundary arc,
and contains the initial nearest prefix.  This is the required pocket;
the prefix crosses each smaller circle once.  After the cap, a stem has
radial depth at least \(\sqrt{15}w\), and the interpolating portion is
at radial distance at least \(20w\) from either stopping layer before
the offset.  Subtracting the offset and curvature/lattice errors leaves
clearance greater than \(w/10\); the retained cap itself has that
clearance by its first-stop definition.  Choose the tube width so that its enlargements are
at most \(w/100\).
Use barrier scale \(\bar R=12R\). All four barriers lie in
\(B_{6R}\) and have radial span at least \(R/8-O(w)>d\), so their
diameters are less than \(\bar R\) and greater than \(d\). Thus
\[
       \theta=\frac{d}{\bar R}
              =\frac{a}{192\,000\,000\pi A}.
\]
The remaining lanes have clearance \(w/10\), so (UB1) holds with
clearance constant \(1/2\) after lattice error. The full endpoint
modules have radius
at most \(d+L\ell_c\le w/100\). Hence the sign-\(\pm\) confined
domains and prescription supports lie in \(\Omega_\pm\) and have
distance at least
\[
             w-w/10-w/100-w/100-O(1)\ge .85w
\]
from \(J\), including lattice error. A free path between opposite
sides cannot cross the fixed \(K\), so it must cross \(J\).
This verifies (UB2) at scale \(\bar R\) with
\(\kappa=.85w/(12R)>0\).
Proposition~\ref{g:barriers} applied to these four signed barriers
forces the interface joins within the disjoint corridors with
uniformly positive probability. Planarity forces the outer arms to
merge with the corresponding inner arms; there is no other exit
through the fences.

The remaining walk bridge must reach the joining vertex of the reversed
inner piece. Its conditioning to reach this vertex before zero has
probability bounded below. Harnack gives a positive probability to
follow the free third corridor to a fixed relative neighborhood of
that vertex; the two-dimensional point-hitting estimate gives
probability comparable to \(1/\log R\) to hit it there. The loop portion
between its first and last visits stays in that corridor with probability
bounded below, by the last-exit decomposition. This proves the lower
bound. Conditional on any completed field configuration, hitting the
same vertex before a nearby interface continuum has probability at
most \(C/\log R\); this proves the upper bound. These are precisely
the walk hypotheses of \cite[Lemma 3.14]{SSdiscrete}: radii in the
specified compact ratio range, observer outside the larger disk,
and two prescribed first/last walk pieces. All field estimates used
in its corridor construction have been supplied above.

We record the two augmented upper bounds used below. First, the upper
bound still holds after the whole field, hence \(\gamma\), is fixed;
the middle walk bridge remains unexposed. Second, put
\(H=(\Phi_{3r},\Theta_{2r})\), and expose that middle walk only until
its first exit from a circle of radius \(\rho/3>10r\), if this occurs
before its prescribed terminal piece. Call this stopped prefix
\(S_{**}\). On this escape event,
\[
 \Prob(Z^\sigma\mid H,S_{**},\hbox{escape})
       \leq C/\log r.                              \tag{GK}
\]
For verification, let \(\hat q\) be the joining vertex of the fixed
reversed inner piece. The first-hit/last-exit decomposition leaves a
walk from the current exit vertex to \(\hat q\), conditioned to hit
\(\hat q\) before zero, followed by its last-exit loop. Harnack bounds
the conditioning probability below by a fixed constant, since the
exit vertex is outside \(B_{10r}\) while
\(|\hat q|=2r+O(1)\). To realize \(Z^\sigma\), this first part
must hit \(\hat q\) before the connected inner interface. The same
point-versus-continuum estimate as above is at most \(C/\log r\);
the remaining loop can only reduce the probability. This argument
holds for every completed field consistent with the records, and
averaging gives (GK). It does not expose the later successful bridge.
\end{proof}

\begin{lemma}[Separation under a harmonic landing]\label{g:separation}
For every \(\eps>0\) there is \(c(\eps)>0\) such that, for all
sufficiently large \(R\),
\[
 \Prob\bigl(Q(\Phi_{3R})\wedge Q(\Theta_{2R})>c(\eps)
       \mid\Phi_{4R},\Theta_R,Z^\sigma\bigr)>1-\eps .       \tag{G2}
\]
The estimate is uniform in admissible records for which the conditioning
has positive probability.  The observer and original boundary are outside
\(\mathfrak B_{12R}\).
\end{lemma}

\begin{proof}
Write \({\cal F}_{r,s}=\sigma(\Phi_r,\Theta_s)\), and write
\(\widehat q_{\rm o},\widehat q_{\rm i}\) for the two numerical quality
formulas above, before the following common live test.  An exact pair is
called open if both pairs of strand tips exist, neither observed interface
has closed or prematurely merged, and the observed walk has made no
forbidden contact with any observed strand or with the original exterior;
the prescribed terminal contact at zero is the sole exception.  The raw
qualities are zero if these conditions fail.  A pair is called live when
\(\Prob(Z^\sigma\mid{\cal F}_{r,s})>0\); its two qualities are the raw
qualities when live and are set to zero otherwise.  A live pair is open.
We shall update only one of the two records at a time.
For the outer record put \(\rho=r\) and \(e=-1\), and for the inner
record put \(\rho=s\) and \(e=1\).  If the selected record has quality
\(q>0\), put
\[
       L=q\rho,\qquad \rho_1=(1+10eq)\rho .
\]
Let \(q_1\) denote its quality after this update, including the live-pair
test, with the other record held fixed.  We first prove, uniformly for
\[
 3.3R<r\le4R,\qquad R\le s<1.7R,\qquad
 0<L<R/200,\qquad L\le\rho_{\rm other}q_{\rm other},
\]
where \(\rho_{\rm other}\) and \(q_{\rm other}\) are the radius and
live quality of the unchanged record,
that
\[
 \Prob\bigl(q_1\ge2q\mid{\cal F}_{r,s}\bigr)\ge c_0>0.
                                                               \tag{G3}
\]
Here \(q=L/\rho<1/200\), so the cap is inactive.  The geometric,
barrier, and walk construction will be proved in the wider layout buffer
\[
 2.8R<r\le4R,\qquad R\le s<1.84R,\qquad
 0<\widehat q\le1/100,
\]
for any jointly admissible open pair, using its selected raw quality
\(\widehat q\), again denoted by \(q\) there.  Only the positive-support
step will use the smaller-
physical-separation premise in (G3).
For the low-quality updates used later we also prove the following joint
form.  Put \(F=\{q_1<2q\}\), and, on \(F\cap\{q_1>0\}\), make one
further update of the same record to
\(\rho_2=(1+10e q_1)\rho_1\), leaving the other record fixed.  Put
\(q_2=0\) on \(\{q_1=0\}\).  If \(q<1/500\), \(r>3.25R\), and
\(s<1.61R\), then
\[
 \Prob(F,q_2=0\mid{\cal F}_{r,s})
       \ge c_1\Prob(F\mid{\cal F}_{r,s}).                    \tag{G4}
\]
Both subsequent radii in (G4) remain in the wider layout buffer just
displayed.
The estimates condition on the stopped records, not on \(Z^\sigma\).

\paragraph{The three local centerlines.}
The selected two strand tips and walk tip have pairwise distance at least
\(L\).  Each strand tip has distance at least \(L\) from the entire old
walk, and the walk tip has distance at least \(L\) from the entire old
strand.  All those histories are on the old side of the rounded circle.
Their deviation from the Euclidean circle is bounded by a lattice constant
\(c_\triangle\).  Below, fixed errors \(O(c_\triangle)\) include the
short access from an actual tip to a circular polar coordinate.  Such
accesses can be chosen in the unobserved side: for an outer tip use the
interior of an incident selected hexagon to reach the Euclidean circle;
for an inner tip use the unselected-hexagon construction given below.

Cut the cyclic order of the three access angles at a largest gap and lift
them to \(a_1<a_2<a_3<a_1+2\pi\).  Write
\(g_i=a_{i+1}-a_i\) for \(i=1,2\) and
\(g_3=2\pi-a_3+a_1\ge2\pi/3\).  Define
\[
 \delta_i=(5q-g_i)_+,\qquad
 (\Delta_1,\Delta_2,\Delta_3)
  =\tfrac12(-\delta_1-\delta_2,\ \delta_1-\delta_2,
                         \ \delta_1+\delta_2).                 \tag{G3a}
\]
Use a fixed tie rule in choosing the cut.  Then \(|\Delta_i|\le5q\).
During interpolation by a parameter \(t\in[0,1]\), the first two gaps
are \(g_i+t\delta_i\); the cut gap is
\(g_3-t(\delta_1+delta_2)\ge2\pi/3-10q>1.99\).
Thus the cyclic order is preserved, no already larger gap is reduced
except the largest one, and all three final circular gaps are at least
\(5q\).

For the moment suppose \(L\ge L_0\), where the fixed threshold \(L_0\)
will absorb all lattice errors.  From the short circular access, retain
a straight radial stem of length \(L/10\).  Spread the angles linearly
to \(a_i+\Delta_i\) by radial progress \(5L\), and keep these angles
fixed for the final radial progress \(5L\), ending at radius \(\rho_1\).
Continue the held rays to total radial progress \(11L\) for the barriers
and to total progress \(13L\) only for a topological completion.  The
spread width is \(4.9L+O(c_\triangle)\).  During spreading the radii
stay between \(.95\rho\) and \(1.05\rho\); using the conservative
factor \(1.1\), after increasing \(L_0\) the ratio of arclength speed
to radial speed is at most
\[
             v_* =\sqrt{1+(5.5/4.89)^2}<1.51.                 \tag{G3b}
\]
The total length of a barrier, including its held extension and the
endpoint construction below, is at most
\(11L+5.5L+O(L/2000)+O(c_\triangle)<17L\).

Here are the required uniform separation estimates.  For \(L_0\) large,
the circular access gaps are at least \(.98q\).  Even through total
progress \(13L\), the common radius is at least \(.87\rho-O(c_\triangle)\).
The corresponding centerline points have distance at least
\(2(.87\rho-O(c_\triangle))\sin(.49q)>.84L\).
If two points have radial difference at least \(L/4\), their distance
has that lower bound.  Otherwise compare at a common radius and use
(G3b), losing less than \(1.51L/4\); the remainder exceeds \(L/4\).
The same argument applies to a stem, a spread piece, or a held extension.
For a short access piece, use its bounded length and the following
old-side estimate.  If \(H\) is an opposite old history and a centerline
point has radial progress \(u\), then
\[
 d(z,H)\ge\max\{L-v_*u-O(c_\triangle),\ u-O(c_\triangle)\}
       \ge {L\over1+v_*}-O(c_\triangle)>.38L.                \tag{G3c}
\]
This also proves the distance from each complete local centerline to the
entire opposite old history.  Consequently any two complete centerlines
are separated by \(L/4\), after one further increase of \(L_0\).

Let \(t_i\) be the three target points at radius \(\rho_1\).  Every
piece of a different centerline before the held part is radially at least
\(5L-O(c_\triangle)\) from \(t_i\).  The distance from \(t_i\) to any
ray whose circular angular distance is at least \(5q\) is at least
\(\rho_1\sin(5q)\): for an angle at least \(\pi/2\) the lower bound
is \(\rho_1\), and otherwise it is the perpendicular distance to the
infinite ray.  This covers the entire held piece and its extension.
Since \(\rho_1\le1.1\rho\) and \(q\le.01\), it follows that
\[
 d(t_i,\hbox{the entire other centerline})
       \ge4.4q\rho_1.                                      \tag{G3d}
\]
This bound includes the bounded access errors.  It is the final held
length \(5L\), rather than only the final angular gap, which controls
the earlier spreading portions.

\paragraph{Four barriers with fixed parameters.}
Put \(w=L/2000\), \(\ell=w/10\), and \(d=w/(200A)\), where \(A\)
is a sufficiently large fixed constant.  Smooth the two centerline
corners in \(O(w)\) neighborhoods, preserving the initial radial stem,
radial monotonicity, and all preceding estimates with a small fixed loss.
The two sides of a regular radius-\(w\) tube about each strand centerline
give its two side lanes.  Such a tube is a strip: the stem and held pieces
are radial, the middle piece is a graph over radius with slope bounded
by (G3b), and the two corners can be rounded with curvature at most
\(1/(10w)\).  The strips for distinct centerlines are disjoint.

For an outer update, complete only the two outer observed arms by these
two centerlines and a simple connector at radial progress \(13L\).
This is a prospective oriented chord \(K_{\rm o}\cup J\).  For an inner
update, use the inner observed strand and the two outward centerlines,
joining at radial progress \(13L\), to obtain a prospective oriented
simple loop.  Let \(K_0\) be the selected old field system, and denote
the two ideal sign sides of \(K_0\cup J\) by \(\Omega_+\) and
\(\Omega_-\).  They are Jordan disks, with the unbounded side of an
inner loop viewed on the sphere.  For an actual attachment in side \(s\),
delete from \(\Omega_s\) the closed lattice triangles meeting \(K_0\), take the component
incident to the connected free collar of \(J\), and take its prime-end
completion.  The deleted triangles attach to the \(K_0\) boundary, so this
component has connected complement and its resolved completion is a disk.
Its non-\(J\) boundary incidences away from tips and marks are the pure
sign-\(s\) bank.  This is the attachment side; a raw pinched component of
\(D(K)\) is not being called a Jordan disk.  The connector is only a
topological separator; no rectangle,
shift, or barrier follows it.  It stays at least \(2L-O(w)\) beyond the
barrier endpoints.  The other stopped field record is disjoint from the
completion because
\[
  (1-13q)r\ge.87(2.8R)>1.84R,\qquad
  (1+13q)s\le1.13(1.84R)<2.8R.                              \tag{G3e}
\]
No compatibility of the sampled inner orientation is assumed here.

We spell out the old-end attachment; it also handles a shadowing return
of the other old arm.  At an actual tip take the circle of radius \(4w\).
The two side lanes meet it at transverse offsets \(\pm w\) and radial
depth \(\sqrt{15}w\), up to \(O(c_\triangle)\).  Follow the circle from
the chosen lane away from the radial stem.  The long arc between these
two lane points stays at distance at least \(w-O(c_\triangle)\) from
the local \(J\); the rest of \(J\) is farther by the stem length and
the centerline separation.  The two lane points are on different sides
of \(K_{\rm o}\cup J\), or of the inner loop.  Thus this arc, which
avoids \(J\), must hit the old field boundary.  At its starting point
the full-boundary clearance exceeds \(3w\): the old record is on the
old side of the rounded circle, and the curvature error is
\(O(w^2/\rho)+O(c_\triangle)=o(w)\).  Stop at its first point \(a\)
of full-boundary distance \(\ell\), and append a globally shortest
segment from \(a\) to \(b\in\partial D(K)\).  The retained circular
arc has clearance at least \(\ell\).  The segment cannot cross \(J\),
since its length is \(\ell<w-O(c_\triangle)\), and it cannot cross a
different fixed boundary before its endpoint.  It is therefore in the
chosen actual side.  Its endpoint lies on an old string, not on the
remote record or exterior, and
\[
 d(b,J)\ge .89w-O(c_\triangle),\qquad
 d(b,\hbox{every observed tip})\ge3.8w-O(c_\triangle).       \tag{G3f}
\]
For later caging, identify this actual endpoint with the ideal bank as
follows.  In the incident deleted triangle selected by its nearest prefix
and meeting \(K_0\), join \(b\) to the first \(K_0\)-hit.  The link has
diameter at most one and stays in the signed ideal side by (G3f).  It is
only an identification through forbidden material, inaccessible to a new
dual continuation away from an active strand tip.  Coincident links are retained as
separate incidences in their boundary order.  This link is not a free path
and is not used to prove the pocket.

The nearest segment has exact clearance
\(d(b+tu,\partial D(K))=t\) for \(0\le t\le\ell\).

This attachment has the pure pocket required by Proposition~\ref{g:barriers}.
Take the radius-\(d\) circular crosscut in the resolved attachment side which
the nearest segment meets.  The circle misses \(J\) by (G3f), so its
two boundary incidences are on the pure bank.  The retained circular arc
has distance at least \(\ell>d\) from \(b\).  From its lane endpoint a
short transverse segment reaches the radial stem at depth about \(4w\);
this segment also has full-boundary clearance greater than \(d\).
It follows that the tail and a collar of \(J\) are on the same side of
the crosscut, while the initial nearest prefix is on the other side.
The crosscut theorem in the Jordan side makes that latter component a
pocket bounded by the crosscut and the single pure boundary subarc which
excludes \(J\).  The nearest prefix crosses every smaller circle once.
The remote record is on the \(J\)-side: in the outer case it lies in
the connected region below radius \(r-L/20\), and in the inner case
in the region above radius \(s+L/20\); these regions reach the stem and
miss the radius-\(d\) crosscut.  Removing the remote record cannot
alter the pocket.  This proof does not assert that the pure subarc or
the pocket has small diameter.

After deleting only the nearest prefix of length \(3d/5\), its remainder
has clearance at least \(3d/5\); the retained circular arc and lanes
have clearance at least \(\ell\), and the free endpoint is remote from
all pins.  Thus (UB1) holds with \(a=1/2\), including the bounded lattice
error.  A point within \(d/5\) of the boundary lies on the nearest prefix
at radius less than \(d/2\), proving the ordinary near-boundary barrier
condition as well.  The endpoint testing window is \(B(b,O(d))\).
Every component of an old simple strand cut by its generous testing
square meets the rim, since all strand tips and the remote record are
outside that square by (G3f).  The incident component also reaches the
rim.  This proves local thickness, even with arbitrarily many returns.
All tested pure pins and the traced pin chain are string pins: these
windows are near radius at most \(4R\), whereas the original exterior
is outside \(12R\).  A remote pure subarc may traverse exterior, but
no exterior pin enters the quantitative attachment estimate; its constant
is the interior-string constant of Proposition~\ref{g:barriers}.

Use the four side barriers only through radial progress \(11L\).
Their radial span exceeds \(d\), and their diameters are less than
\(17L\).  Take common barrier scale \(\bar R=32L\).  Choose the fixed
mesh and interpolation neighborhoods in the proposition so that every
full endpoint module has radius at most \(w/100\) and every tube
enlargement has width at most \(w/100\).  The full domains \(C_\pm\)
and all supports \(P_\pm\) lie in their corresponding \(\Omega_\pm\)
at distance at least \(.8w\) from \(J\), by the lane offset, (G3f), and
the two \(w/100\) losses.  A free path between opposite sides cannot
traverse fixed \(K\), hence must cross \(J\).  This proves (UB2).
For the ordinary internal distance, an access-labelled admissible path
has its interior in \(U\).  Under the preceding forbidden-material bank
identifications it starts on opposite ideal sides of the continuous
separator \(K_0\cup J\); it cannot meet \(K_0\), and therefore crosses
\(J\cap U\).  Its diameter is at least \(.8w>2d\).  Thus the ordinary
opposite-barrier separation also holds, with the fixed choices
\[
 \theta={d\over32L}={1\over12\,800\,000A},\qquad
 m=4,\quad a=1/2,\quad \kappa={.8w\over32L}={1\over80\,000}.   \tag{G3g}
\]
The observed component near an endpoint supplies the required anchor.
There are a bounded number of smooth pieces and turns, and all rectangle
counts in the proposition depend only on (G3g), not on \(q\).  The
probability of avoiding all four barriers is therefore at least a fixed
\(p_{\rm b}>0\) under either actual inner orientation.

\paragraph{What the barriers force.}
Add only for topology a transverse gate from each side barrier's free
endpoint to the corresponding centerline at radial progress \(11L\).
The two gates for tip \(i\) end at the same point \(g_i\in J\), beyond
the stopping circle.  In the ideal \(\Omega_s\), adjoin the forbidden-
material endpoint link; the side barrier plus this link and gate is a proper
crosscut \(A_{i,s}\) from the pure bank \(B_s\) to \(J\).
The two same-sign crosscuts are disjoint.  Parametrize \(B_s\) and \(J\)
from tip 1 to tip 2.  Their \(J\)-endpoints have this order, so their
\(B_s\)-endpoints have the same order: otherwise the four endpoints
alternate and the crosscuts must meet.  This conclusion permits both
attachments to be on shadowing parts of the old strand.

Glue the two Jordan sides along \(J\), retaining distinct incidences
on the two fixed banks.  The union \(Y_i=A_{i,+}\cup A_{i,-}\) is a
proper crosscut in this split-bank domain.  The preceding nonalternation
gives disjoint tip-1 and tip-2 end components and a central component.
Extra fixed observations can only remove points.  In the relative boundary
of an end component the two virtual \(J\) portions cancel; its only exits
are actual barriers, fixed boundary, and its beyond-stop gates.  A fresh
interface prefix which avoids the barriers and is stopped at the new
circle therefore stays in its end component.  The two prefixes cannot
merge there.  For an outer update the completed chord must connect the two
old tips; for an inner update the complete component either joins the two
continuations or reaches the original exterior, which is beyond the new
circle.  Thus neither case can lose a target tip before the stopping circle.
On the new circle all cap and turning pieces are behind,
and each held lane crosses in one bounded lattice block.  The end component
meets that circle precisely in the short interval between its two lane
blocks.  Each new strand tip is consequently within \(O(w+c_\triangle)\)
of its target.  No metric confinement of the rest of that prefix in a
narrow tube is asserted.

Put \(\nu=L/1000\).  Require the fresh walk to follow the width-\(\nu\)
tube about its centerline through a terminal extension of length \(2\nu\),
including a full entrance ball of radius \(L/100\) about its starting
vertex and the radius-\(\nu\) terminal ball.  The terminal ball extends
at least \(\nu\) beyond the new circle.  The entire corridor meets that
circle only within \(\nu+O(c_\triangle)\) of the target: the entrance ball
is \(10L\) away, the spreading pieces are at least \(5L\) away, and the
remaining tube follows the held ray.  Stop the successful walk at its first
crossing of the new rounded circle in this window.  The
entrance ball is not clipped at the old circle and may meet the old walk.
By (G3c) it misses the old field, and centerline separation makes the
whole tube and ball miss the side barriers and their gates, including the
hexagon error.  They lie in \(B_{5R}\), so they also miss the original
exterior outside \(\mathfrak B_{12R}\).  Let \(X\) be the ideal ambient
disk or sphere slit only along the selected old dual system \(K_0\), with
its bank accesses retained and the short forbidden endpoint links included
in \(Y_1\cup Y_2\).  Let \(W\) be the union of the closed hexagons of
every old and fresh visited walk vertex, not just the polygonal walk edges.
Every cell interior avoids \(K_0\), which lies on hexagon edges.
Consecutive visited cells share an edge not crossed by \(K_0\): for the
old record this is the open no-prior-contact condition, including the
reversed first step from zero, since otherwise its forward predecessor
would already be adjacent to the interface; for the fresh part it is the
tube clearance.  Thus \(W\setminus K_0\) has a connected lift to \(X\),
with the prescribed access along the old contacts of the origin cell.
Those boundary contacts are not asserted to be in the open surface.
The old walk misses the crosscuts because away from the \(O(w)\) tip caps
they lie on the new side of the old circle, while the old quality protects
the caps.  Together with the tube separation, this makes \(W\) disjoint
from \(Y_1\cup Y_2\), including their links and gates, after the hexagon
error.  The terminal cell meets the central circle interval.  Hence the
connected lift of \(W\setminus K_0\) lies in the central component of
\(X\setminus(Y_1\cup Y_2)\), even if it crosses virtual \(J\).  A fresh
dual prefix stays in an end component, so it cannot meet \(W\) off
\(K_0\).  It cannot newly enter \(K_0\) away from an active strand tip; the
pre-existing origin contact remains allowed.

Put a ball about the deterministic walk target of radius
\(2q\rho_1+\nu+O(c_\triangle)\), including the full landing-window
and hexagon errors.  Equations (G3c)--(G3d), the \(5L\) radial hold,
and the \(O(w)\) gates show that this ball misses both paired crosscuts
\(Y_i\).  It also misses the old field by the \(10L\) radial gap.
It misses the remote stopped record by (G3e), including this ball radius.
The ball is connected and contains the central target; its lift to the
split-bank domain is wholly central even if it crosses virtual \(J\).
Thus the actual walk tip is at distance at least \(2q\rho_1\) from the
entire newly observed field.  Conversely each new field tip is within
\(O(w+c_\triangle)\) of its target; (G3d) bounds its distance to the
entire new walk tube, and the radial gap bounds its distance to the entire
old walk.  Both exceed \(2q\rho_1\).  Finally the two field targets have
distance at least \(2\rho_1\sin(5q/2)\), which still exceeds
\(2q\rho_1\) after the landing errors.  Hence every term in the new
quality numerator is at least \(2q\rho_1\), and no forbidden contact
or premature merger has occurred.

\paragraph{The walk probability.}
For \(x,y\ne0\), put
\[
 G_0(x,y)=\sum_{n\ge0}\Prob_x(S_n=y,\ n<T_0).
\]
This killed Green function is finite and symmetric.  If the exact outer
and reversed inner walk records end at \(a\) and \(b\), the unexposed
middle path has law
\[
 \mathsf B_{a,b}(m)={6^{-|m|}\over G_0(a,b)},\qquad
 m:a\longrightarrow b,\quad m\cap\{0\}=\varnothing.          \tag{GW1}
\]
The fixed prefix and suffix contribute a common factor to the walk weight;
summing the remaining weights gives \(G_0(a,b)\).  All revisits of either
recorded history or circle are permitted.  Reversal is exactly
\(\mathsf B_{b,a}\), including all loops, and the field record does not
change this law before conditioning on \(Z^\sigma\).

Let \(E\) be the successful entire-corridor first-crossing event specified
above, stopped at \(\tau\).  Summing the still unexposed remainder in
(GW1) gives for an outer extension
\[
 \mathsf B_{a,b}(E)
       ={\E_a[\mathbf1_EG_0(S_\tau,b)]\over G_0(a,b)}.       \tag{GW2}
\]
The full entrance ball and width-\(\nu\) tube have length \(O(L)\) and
are covered by a fixed number of overlapping balls of radius comparable
to \(L\).  Their entire corridor and a fixed relative enlargement containing
the doubled Harnack balls stay between radii \(.89r\) and \(1.01r\) for
an outer update, or between \(.99s\) and \(1.12s\) for an inner update.
The remote-pole gaps are bounded below by
\[
 .89(2.8R)-1.84R>.65R,\qquad
 2.8R-1.12(1.84R)>.73R,
\]
and the enlargement is also a fixed positive distance from zero.  Since
\(L\le.04R\), these are uniform relative Harnack margins.  Ordinary
walk traverses the ball chain with a fixed positive probability.  Positive
harmonicity of \(G_0(\cdot,b)\) on its doubled balls compares its value
on the landing window with \(G_0(a,b)\), so (GW2) is at least a fixed
\(p_{\rm w}>0\).  The full entrance ball permits the necessary old-circle
recrossings.  After any exact outer prefix ending at \(a'\), the remaining
law is \(\mathsf B_{a',b}\); reversing it gives the inner estimate with
\(G_0(\cdot,a')\).  The current stopped layers remain strictly nested,
and the same displayed pole margins hold for every current joining vertex
in the buffer.  Thus this also applies after an earlier one-record update,
without an independence assertion.  Field and walk are independent given
their respective stopped records and the actual field orientation.  The
bound \(p_{\rm b}p_{\rm w}\) survives averaging over that orientation,
and gives an open raw output of quality at least
\((2q)\wedge1/100\) with this uniform probability for \(L\ge L_0\).
This part has not used the live test.

\paragraph{Bounded lattice scales.}
We use a simultaneous form of bounded insertion.  Given a fixed number
\(m\) of fixed-size patches near observed tips and compatible sign
prescriptions on a bounded number of their free vertices, their joint
probability under the original exact-strand conditioning has a positive
uniform lower bound.  Fix first the shell failure parameter \(1/(4m)\).
For each fixed patch size choose the corresponding long-slack enlargement
from the proof of Lemma~\ref{g:insertion}, with this parameter and its
required deficit.  Merge clusters whenever their already chosen enlarged
boxes, with one neighbor layer, intersect, recomputing the enclosing box
with the same parameter after a merge.  There are at most \(m-1\) merges,
and overlap bounds each new diameter by a deterministic function of the
previous ones.  Thus the final boxes have uniformly bounded sizes and no
edges between their interiors; there is no later enlargement.  The arbitrary
animal parameter in Corollary~\ref{f:conditioned-sparse} makes each shell
failure less than \(1/(4m)\).  With probability at least \(3/4\) all
shells are good.  Construct the prescribed feasible filling separately
inside every box, keeping all shells fixed.  On their disjoint interiors
perform one joint contraction toward these fillings.  Its Jacobian is
\(t^{\sum N_i}\), with \(t>0\) and \(\sum N_i\) bounded by the patch
parameters.  Integrating proves the joint bound.  No first fence pattern
is used as a new conditioning for a second invocation.

For \(L\le20\) use the deterministic triangular/hexagonal routing in
\cite[proof of Lemma~3.15, pp.~65--66]{SSdiscrete}.  It uses only the
local rounded-circle geometry, the full-history quality inequalities,
and positive completion support.  Its longest boundary arc and extremal
left/right dual continuations leave a bounded primal route avoiding all
adjacent hexagons, with doubled quality and preserved completion support.
The same construction with first exit in place of first entry gives the
inner route; the fixed ratio \(r/s>1.52\) keeps the opposite record out
of its bounded patches.  The simultaneous insertion argument just proved
replaces only the field probability input.  A prescribed primal route of
at most \(M\) steps has, under either Green transform in (GW2), probability
at least \(6^{-2M}\): for the current remote pole \(b\), harmonicity gives
\(G_0(v_{j+1},b)\ge G_0(v_j,b)/6\) at each step.  Thus this branch has a uniform
positive probability without drawing sublattice-width barriers.

For \(20<L<L_0\) we give a direct finite construction.  Put
\(h_{\mathrm{hex}}=1/\sqrt3\).  For an outer strand tip, join it through an
incident selected hexagon to a point of the Euclidean radius-\(r\) circle;
this costs at most \(2h_{\mathrm{hex}}\) and avoids the old dual strand.  For
an inner tip choose an incident unselected hexagon with center \(v\),
join the tip to \(v\) through its interior, and choose a lattice unit
vector \(e_v\) with \(e_v\cdot v/|v|\ge\sqrt3/2\).  An unselected
hexagon has \(|v|\ge s+1/2\), since it contains the radius-\(1/2\) disk
about its center, and boundary incidence gives \(|v|\le s+3h_{\mathrm{hex}}\).
For every integer \(k\ge1\),
\[
 |v+ke_v|\ge |v|+k\sqrt3/2>s+h_{\mathrm{hex}},
\]
so these successive hexagons are unselected.  Their center segments
remain in the unselected union.  Follow them until the radius-\(s+7/4\)
circle is first met.  This occurs after center-path length at most
\((7/4-1/2)/(\sqrt3/2)<1.444\); the whole tip access has length less
than \(2.022\) and avoids the inner strand.  Use the radial circular
access for the walk, with its bounded first-entry or first-exit error.

The outer access points have pair distance at least
\(L-4/\sqrt3>.88L\).  Use (G3a) with its signed shifts bounded by
\(4.12q\), an initial radial prefix to radius \(r-8\), spreading to
\(r-5L\), and the final held piece to \(r-10L\).  The spread width is
at least \(4.6L\).  Comparing spread parameters at threshold \(1/12\)
gives either radial separation \(4.6L/12>.383L\), or angular gap at
least \((.88-4.12/12)q>.536q\), whose chord exceeds \(.50L\).
The same comparison handles a prefix against a spread or held piece.
The short access pieces have length at most \(2/\sqrt3\) and stay within
a bounded vanishing radial sagitta of the old circle.  Against another
prefix, either their radial gap is at least \(7.7\), or that prefix is
within \(2/\sqrt3+7.7+o(1)\) of its own tip; the original separation then
leaves more than \(20-4/\sqrt3-7.7-o(1)>9.9\).  Against a spread piece
the radial gap is at least \(8-o(1)>7.66\), and held pieces are farther.
Thus the outer complete centerlines have mutual distance \(>7.66\).
Their speed ratio is at most \(\sqrt{1+(4.12/4.6)^2}<1.35\), so the
outer analogue of (G3c), with old histories outside radius \(r-1\), is
\[
 d(z,H)\ge\max\{L-2/\sqrt3-1.35u,\ u-1\}>.36L>7.2.       \tag{G3i}
\]

For the inner access circle the pair distance is at least
\(L-4.044>.797L\); since \(s\ge2000\), its angular gaps are at least
\(.796q\).  The centered shifts in (G3a) are at most \(4.204q\).
Take the straight prefix to radius \(s+10\), spread to \(s+5L\), and
hold to \(s+10L\).  The spread width is at least \(4.5L\).  At parameter
threshold \(1/11\), the two alternatives give radial separation
\(4.5L/11>.409L\) or angular gap
\((.796-4.204/11)q>.413q\), whose chord exceeds \(.412L\).
A short access has radius at most \(s+7/4\), whereas a spread piece has
radius at least \(s+10\), giving gap \(8.25\).  Against another prefix,
either the radial gap is at least \(7.8\), or that prefix is within
\(2.022+7.8\) of its own tip; the original tip separation then leaves
more than \(20-2.022-(2.022+7.8)>8.15\).  Hence all inner complete
centerlines have mutual distance greater than \(7.8\).  The old inner
strand is inside radius \(s+2/\sqrt3\) and the whole old reversed walk
inside radius \(s+1+1/\sqrt3\).  The speed ratio after the access is
less than \(1.41\).  The counterpart of (G3c) is therefore greater
than \(7.5\) for \(L\ge20\); explicitly, for radial depth \(u\ge7/4\),
it is at least
\[
 \max\{L-2.022-1.41(u-7/4),\ u-1-1/\sqrt3\}>7.5.          \tag{G3j}
\]
The short access itself has distance at least \(L-2.022\) from the
opposite old history.

Use finite side lanes of width \(w_b=3\) and first full-boundary stop
distance \(\ell_b=4/3\).  Near a tip take the appropriate one-sided
outer boundary of the radius-three neighborhood of its short access and
straight prefix, truncated at outer depth six or inner depth eight.  The
short access has diameter less than three, so its radius-three disks have
a common interior core; adjoining the monotone radial prefix gives a local
strip with an end cap.  Its side frontier has distance three from the local
\(J\).  The future spread starts two radial units farther.  Its tangential
slope is at most \(.982\), so even in the direction of the offset point its
distance there is at least
\(\min_{t\ge0}\sqrt{(2+t)^2+(3-.982t)_+^2}>3.5\), up to a vanishing
polar-curvature error; earlier points of the end collar are farther.
Other centerlines are farther than \(7.66\).  Thus this local frontier
has distance three from the full \(J\), and it joins the side boundary of
the full radius-three tube.  Its deep point has full deleted-boundary
clearance greater than \(6-1-o(1)>4\) in the outer case and
\(8-2/\sqrt3-1-o(1)>5.8\) in the inner case.

The old strand must meet this frontier near the tip.  For the outer case take its first point at
distance ten from the tip.  A \(J\)-point of inward depth at least four
is at least four from the old strand; a shallower one is within
\(4+2/\sqrt3<5.16\) of its own tip, giving distance greater than
\(10-5.16>3\) from that old point, for either local \(J\).
For the inner case a \(J\)-point of outward depth at least five is at
least \(5-2/\sqrt3>3\) from the old strand; a shallower one is within
\(2.022+(5-7/4)<5.273\) of its own tip, again leaving distance greater
than \(10-5.273>3\).  These old points exist because the relevant old
strand continues a macroscopic distance to the exterior or to the
incident-edge neighborhood.  Its arc starts on \(J\) and exits the
radius-three neighborhood.  It cannot exit through the deep cross-section,
which is strictly on the new side of the old radial bound.  In each actual
Jordan side the local strip frontier from the deep lane must therefore end
on the old boundary; an earlier shadowing return only ends it sooner.
Stop at the first clearance
\(\ell_b\) and nearest-attach.  Every point of this open finite barrier
then has distance at least
\[
                  3-4/3=5/3>2/\sqrt3                       \tag{G3k}
\]
from \(J\).  Continue along the held lane to a free endpoint beyond the
new circle and add the same virtual gate as above.

To identify the actual endpoint \(b\in\partial D(K)\) with its bank of
the central chord or loop, choose the incident deleted triangle which
meets \(K\), and join \(b\) inside that triangle to its first \(K\)-hit.
The triangle has diameter one, less than \(d(b,J)\ge5/3\), so this link
stays on the same bank.  It is only an identification inside forbidden
material, not a free path or a sign prescription.  Coincident links are
resolved in their boundary order.  An actual dual continuation cannot use
such a link or enter its endpoint away from an observed tip.

Perturb these finitely long barriers off dual vertices.  Prescribe sign
\(s\) at every free primal vertex whose closed hexagon meets an open
sign-\(s\) barrier.  An already fixed hexagon meeting it has sign \(s\):
its interior lies on that bank of the prospective chord or loop, and by
(G3k) it cannot also meet \(J\).  Remote fixed observations miss these
bounded lanes.  No free hexagon meets opposite-sign barriers, for a path
through that hexagon between the two sides would have to cross \(J\),
requiring diameter at least \(5/3\), greater than the hexagon diameter.
These are compatible prescriptions on at most \(M(L_0)\) vertices in
at most two bounded patches near the old tips.  The simultaneous insertion
argument gives them a joint positive probability.  At a transverse crossing
of an open barrier, both hexagons adjacent to the crossed dual edge meet
that barrier and have its strict sign.  A zero interface cannot cross it,
nor pass through its fixed endpoint away from a tip.  The paired-crosscut
caging proof above therefore applies literally.

For the walk, the cells successively met by its centerline give a primal
nearest-neighbor route in its full radius-one tube; at a cell-vertex tie
insert an incident neighboring cell.  Include the full radius-one entrance
ball and extend the route two units beyond the target circle.  In the inner
case the final crossed cell center has radius at least
\(s+10L+2-1/\sqrt3>s+10L+1/\sqrt3\), so it is unselected and the route
has exited.  Stop at first exit; use first entry in the outer case.  The
route has at most \(M(L_0)\) steps and controls the whole new walk history,
including revisits.  The fixed radial buffer keeps every route vertex away
from zero and stops it before the remote joining vertex.  Equations
(G3i)--(G3j) and the centerline separations
leave fixed lattice room between this tube and the finite sign fences.
The new field tips lie within five lattice units of their targets.  The
targets are at radial distance \(5L\) from all earlier pieces and at
distance at least \(4.8q\rho_1\) from every other held ray; subtracting
the target, lane, and hexagon errors (at most eight units) still leaves
more than \(2q\rho_1\) for \(L>20\).  Old histories are separated by
the \(10L-O(1)\) radial gap.  The expanded target-ball and central-
component argument controls the entire fresh field, as before.  The
finite route probability is at least \(6^{-2M(L_0)}\) by (GW2).
Thus the direct construction gives an open raw output of quality at least
\((2q)\wedge1/100\) with a uniform positive probability on the remaining
bounded range, before the live test.  The same deterministic construction
works with any fixed upper bound in place of \(L_0\); only the finite
cell count and its probability constant change.

\paragraph{Positive support after direct steering.}
The branch \(L\le20\) already includes the completion support in its
cited discrete routing conclusion.  Suppose now that \(L>20\) in (G3),
and fix any exact output pair in the successful open raw event just
constructed.  The unchanged physical separation is at least \(L\), by
the hypothesis of (G3), and the updated one is at least
\(2(1-10q)L>1.9L\).  Thus both raw physical separations exceed twenty.
The input live pair has a compatible cyclic matching of its two strand
accesses and its walk access, obtained from any completion witnessing
\(Z^\sigma\).  The target order and the landing error bounds show that
direct steering preserves this matching.  It is a geometric matching;
the inner orientation sampled for the raw event is not retained.

Fix a sufficiently large absolute \(H\), to be specified below.  From
the fixed output pair prospectively update any raw physical separation
\(\ell<H\), leaving the other record fixed, until both are at least
\(H\).  Use the direct finite sign-fence construction above with its
fixed upper range enlarged to \(\max\{L_0,H\}\).  All its deterministic
inequalities hold for \(L>20\); only its finite cell count changes with
this upper range.  This construction and its finite walk route use only
joint admissibility and openness, not the live test.  In particular they
give, with positive conditional probability, an open output with raw
quality at least twice the input raw quality whenever the cap is inactive.
The fresh walk remains inside \(B_{5R}\), so it creates no contact with
the original exterior.  After each such event forget its auxiliary sign
prescriptions and orientation and retain an admissible exact output pair.
There are countably many exact path outputs, so one has positive
conditional probability.  This is a prospective positive-support
construction, not a uniform lower bound.

For \(R\ge1000H\), every selected raw quality in this continuation is
at most \(H/R\le10^{-3}\).  Each successful update therefore multiplies
its physical separation by at least \(1.8\).  For either record the sum
of its own starting physical separations before reaching \(H\) is less
than \(1.8H/(1.8-1)=2.25H\); its total radial change is less than
\(22.5H\).  The original (G3) output has \(r>3.135R\) and
\(s<1.785R\).  Hence the whole prospective continuation remains in the
wider \(2.8R,1.84R\) layout buffer.  It ends after finitely many updates,
and retains openness, the cyclic matching, and the absence of every
forbidden strand or exterior contact.

At this final raw pair use only the deterministic centerlines and capped
side curves from the proof of Lemma~\ref{g:hookup}, not its quantitative
probability assertion.  Take \(a=H/(4R)\).  Both normalized raw qualities
are at least \(a\); the rounded-circle errors are absorbed by choosing
\(H\) large.  The deterministic estimates there give centerline
separation at least \(cH\) and lane width at least \(w_H=c'H\), for
fixed absolute \(c,c'>0\).  Fix \(H\) so large that \(w_H\) exceeds one
hundred times every absolute cell error below.  The compatible cyclic
matching chooses the orientation of the inner strand and the side
\(s_0\) containing both walk accesses.  Write \(\Gamma=K\cup J\)
for the resulting simple ideal chord, with its two sign sides
\(\Omega_\pm\).  Here \(J\) consists of the two free connectors in
the unobserved annulus, with the radial end stems just specified in that
proof.

Use the sign-\(+\) and sign-\(-\) side curves along just one of these
connectors.  At an endpoint concatenate the open curve with a link through
the selected same-side sector of its incident deleted triangle, ending
in the interior of a pure strand-adjacent hexagon.  If necessary the link
uses at most two old hexagons of that same sign.  It does not include the
actual \(K\) endpoint: at a tiling vertex use the specified access and
end in a pure edge-side cell.  Perturb the concatenated curve within its
side so that it crosses tiling edges transversely and misses tiling
vertices.  The list of crossed hexagons is then an edge-connected primal
chain from the outer sign bank to the inner sign bank.  Every old sign
site on this chain has the indicated sign, because its hexagon interior
lies in \(\Omega_\pm\).  No hexagon belongs to both chains: its interior
avoids \(K\), and a path inside it between the two sides would have to
cross \(J\), contradicting the distance of each open side curve and
endpoint link from \(J\), which is greater than the hexagon diameter.

Join the two walk accesses by the third deterministic centerline,
including its actual walk-tip access pieces, and list the hexagons it
crosses after a generic perturbation.  Their centers give a finite primal
route between the exact walk-tip vertices; its full neighbor layer stays
in the free third corridor by the choice of \(H\).  Concatenate it with
the outer record and the reversal of the inner record to obtain a
candidate path \(\Pi\) ending at zero.  For
every nonzero vertex of \(\Pi\), request sign \(s_0\) on that vertex
and all its lattice neighbors.  These requests are compatible with all
old signs.  Indeed the old walked hexagons cannot cross \(J\): their
histories are on the old sides of the stopping circles, and the only
end neighborhoods of \(J\) are separated from those histories by the
physical gap \(H\), including one neighbor layer.  They cannot cross
\(K\) by openness.  Their accesses and the middle route are in
\(\Omega_{s_0}\), so all their nonzero hexagons lie on that side.
If a neighboring old sign hexagon were on the opposite side, the common
primal edge would meet \(\Gamma\).  A \(K\) edge there is exactly a
forbidden prior adjacency, and a \(J\) crossing is excluded by the
preceding clearance.  An exterior neighbor on the opposite sign arc
would likewise make the common primal edge cross \(K\) or \(J\),
contrary to observed no-contact or the same clearance; a same-side
Dirichlet neighbor is retained unchanged.  The same
side and hexagon-diameter argument excludes an overlap with the opposite
sign chain.  Requests of sign \(s_0\) cause no conflict with the other
chain.  The origin may occur in a neighbor layer and then has the
compatible sign \(s_0\); only the star centered at the origin is omitted.

Choose \(0<b<1/4\).  Give every requested free sign site, including all
old string sites, the value \(+b\) or \(-b\), give all other free sites
value zero, and retain the original Dirichlet pins.  No site has received
opposite requests.  The retained old signs include the literal tail
signs through each crossing triangle and first edge on the other side of
its rounded stopping layer; these preserve the specified first entrance
or exit and the exact in-disk component.  Each free--free edge has height difference at most
\(2b<1\), and each free--pin edge has difference at most
\(b+1/2<1\); the pin--pin constraints were already feasible.  A
relative-open neighborhood therefore satisfies all old exact strand
signs, the two chains, the nonzero walk stars, and the hard Lipschitz
constraints.  It has positive volume in this finite domain, and in
particular proves positive probability for the chosen inner orientation.

Along a simple zero strand the consecutive same-side hexagons are equal
or edge-adjacent.  Each of the two forced chains consequently joins the
corresponding sign next to \(e_\sigma^*\), through the inner bank and
outer bank, to its corresponding exterior sign arc.  The zero component
through \(e_\sigma^*\) cannot be a closed cycle, since a cycle would
separate one of these two sign clusters from the exterior.  Interior zero
vertices have degree two and the only boundary sign changes are the two
marks, so this component is the boundary chord.  At every nonzero
vertex of \(\Pi\), all neighbors have its one strict sign; no zero edge
is incident to its hexagon.  Thus \(\Pi\) meets the chord first at its
allowed endpoint zero.  It has no earlier exterior contact by the same
raw openness and corridor clearance.  The finite middle walk route has
positive conditional Green-bridge probability.  The strict field event
and this route therefore give \(Z^\sigma\) positive conditional
probability at the terminal \(H\)-pair.  Each of the finitely many raw
success transitions leading there had positive conditional probability,
with its auxiliary orientation and fences forgotten at the retained exact
record.  Applying the tower property successively along this finite
continuation gives positive hookup probability at the original fixed
direct output pair.  Every direct successful output in (G3) is live.
Combining this with the uniform raw steering bound
and the already supported \(L\le20\) branch proves (G3).

\paragraph{The one-record collapse estimate.}
For the outer update the proof of (G4) follows the stopped first-approach
decomposition in \cite[proof of (3.38)]{SSdiscrete}; we record the exact
conditioning because it is needed below.  Reveal the new two field arms
first, but not the whole intervening walk.  Write \(x,y\) for their tips,
\(z\) for the eventual new walk tip, and \(S^\dagger\) for the newly
exposed walk portion.  The failure event \(F\) is covered, including ties,
by the following events:
\[
\begin{split}
 M_0&=\{q_1=0\},\\
 M_1&=\{d(S^\dagger,\{x,y\})<(2q\rho_1)\wedge|x-y|\},\\
 M_2&=F\cap\{\rho_1q_1=d(z,\hbox{new whole strand history})\},\\
 M_3&=F\cap\{0<\rho_1q_1=|x-y|\}.
\end{split}                                                   \tag{G4a}
\]
Events involving missing tips are empty, that case being in \(M_0\).
The cap is inactive because \(2q<1/100\).  The old opposite histories
cannot make the new quality insufficient: their radial gap is \(10L\),
whereas \(2q\rho_1\le2.2L\).  At the bounded lattice scales this is
understood with the exact first-entry/first-exit version of the cited
discrete construction.

For \(M_1\), stop \(S^\dagger\) at its first visit within
\(a=(2q\rho_1)\wedge|x-y|\) of either new strand tip.  Condition only
on the new field arms and this stopped walk prefix, and call the approached
tip \(x\).  The other tip is at least \(a\) away.  The preceding annular
width is \(10q\rho\), whereas \(a\le2q\rho_1\); for either direction
their ratio is at least \(5/(1+10q)>250/51\).  At diverging scales the
rounded-circle error therefore still leaves width at least \(4.8a\).
The approached strand crosses the annulus about \(x\) from radius
\(a/4\) to radius \(2a\).  Starting at this first approach, stop a local
trial at a field hit, entry to either radius-\(a/4\) tip ball, or exit
from \(B(x,2a)\).  The source's planar stopped-walk estimate gives a
fixed probability of the field-hit alternative.  Its fixed relative
enlargement avoids the remote joining vertex and zero in the present
buffer, so Harnack and (GW2) give the same assertion for the Green bridge.

The trial need not finish before the first crossing of the
\(\rho_1\)-circle.  If that crossing occurs first at \(z\) with
\(d(z,H)<a/4\), where \(H\) is the whole new strand history, interrupt
the trial there.  Before that time the walk has stayed outside the two
radius-\(a/4\) tip balls, and before its first approach it stayed at
distance at least \(a\) from both tips.  Since \(|x-y|\ge a\) and the
cap is larger than \(a/4\), either the pair is already dead or
\(\rho_1q_1=d(z,H)\).  The direct \(M_2\) hitting estimate at this
complete first record, with scale \(d(z,H)\), then kills hookup during
the next update with a fixed probability.  If instead the crossing has
\(d(z,H)\ge a/4\), every quality term at that crossing is at least
\(a/4\), unless hookup is already dead.  Its next radial displacement
is therefore at least \(10(a/4)=2.5a\).  A field hit in the local trial,
which stays in \(B(x,2a)\), occurs before that next circle; the strict
\(0.5a\) margin absorbs the rounded-circle error at diverging scales.
A trial hit before the first crossing also kills hookup.  Thus the union
of a local hit before interruption and an interruption contains the
un-interrupted local-hit event, so it has a fixed lower probability.  On
the interruption event, condition on the record at interruption, not on
whether the trial would later hit; the direct estimate then supplies a
fixed conditional death probability.  Conditioning there and integrating
proves
\[
 \Prob(q_2=0,M_1\mid{\cal F}_{r,s})
       \ge c\Prob(M_1\mid{\cal F}_{r,s}).                    \tag{G4b}
\]
This is not asserted after fixing the later escaping part of the failed
walk.  On \(M_2\) the same direct hitting estimate starts at the new
walk tip at scale \(\rho_1q_1\), and may be used after that full new
record is fixed.  Its next radial displacement is ten times this scale.
At bounded scales the exact discrete first-approach and finite routing
argument of the source, with bounded insertion, gives the same fixed
lower bounds.  On \(M_3\), one
single-sign circular-clearance barrier around the midpoint of \([x,y]\)
forces the strands to join before the next circle, whose radial gap is
\(10|x-y|\).  Its nearest endpoints, pure pockets, and local thickness
are those proved above for the circular construction, and the remote record
is outside its annulus.  Proposition~\ref{g:barriers} applies at diverging
scales and bounded insertion at bounded scales.  \(M_0\) is immediate.
Thus (G4b) holds for all four \(M_i\) with one \(c>0\).  Summing the
possibly overlapping events gives (G4) with \(c_1=c/4\).

For an inner update use the same decomposition for the reversed first-exit
walk and the outward-growing strand.  The entire old histories now lie
inside the old rounded circle, so the radial-gap calculation reverses
without changing its constants.  Reverse the whole Green bridge as in
(GW1)--(GW2); its local first-approach law has the remote outer joining vertex
as pole.  The same first-approach hitting argument and the reflected
single-sign circular barrier apply.  The quantitative buffers ensure
these local regions miss the fixed outer record.  This proves the joint
form (G4) for either selected record, without claiming a hazard bound
conditional on a complete failed intermediate record.
We now iterate the one-record estimates at the smaller physical
separation.  Write
\[
 \ell_{\rm o}=r q_{\rm o},\qquad \ell_{\rm i}=s q_{\rm i},\qquad
 p={\ell_{\rm o}\wedge\ell_{\rm i}\over R},\qquad \beta={3\over2},
\]
where both live qualities are zero at death.  Fix
\(a_0=\beta^{-N}\le10^{-3}\), to be chosen below, and put
\[
 {\cal B}=\{(r,s):7R/2\le r\le4R,\ R\le s\le3R/2\}.
\]
Starting from \((r,s)=(4R,R)\), stop at death, when \(p\ge a_0\), or
when the radii leave \({\cal B}\).  Stopping is checked only at round
boundaries; let \(\tau\) be the first terminal boundary.

At a live round boundary choose a record with smaller physical separation
\(\ell=pR\), using a deterministic tie rule.  Its quality is
\(q=\ell/\rho\le p<a_0\), so the cap in (G3) is inactive.  Attempt its
(G3) update, leaving the other record fixed.  On success its new physical
separation satisfies
\[
 \ell_1=\rho_1q_1\ge2(1+10eq)\ell\ge1.8\ell.              \tag{G5a}
\]
On failure \(q_1<2q\), flag the round and immediately execute the same
record's (G4) update.  If \(q_1=0\), end dead without another radial
update.  Otherwise use \(q_1\) to set the collapse radius and then end
the round.  After a first success, let \(U\ge\ell\) be the unchanged
physical separation.  If \(U\ge\beta\ell\), end the round.  If
\(U<\beta\ell\), it is now the smaller separation by (G5a); attempt its
(G3) update.  Its quality is at most \(U/R<\beta p<.0015\), so its cap
is also inactive and (G4) applies on failure.  Flag and immediately
collapse that same record on failure, again making no further update if
its failed quality is zero.  End the round after this second attempt or
its collapse.  A flag-free round has \(p_{k+1}\ge\beta p_k\).

A failed growth and its following collapse are one transition for the
round filtration: the intervening full record is not added before (G4)
is used.  The final nested records determine that restriction, the chosen
radii, and the flag, without observing any still unexposed bridge.  Let
\(\mathcal F_k\) be the exact pair at the start of round \(k\), let
\(J_k\) indicate its flag, and let \(A_k\) mean that the state is live.
Decrease \(c_1\), if necessary, so \(0<c_1<1\), and put
\(\varrho=1-c_1\).  The joint form of (G4) says that for either failed
attempt \(F\),
\[
 \Prob(F,\text{ survival after the immediate collapse}\mid
                  \mathcal F_{\rm pre})
 \le\varrho\Prob(F\mid\mathcal F_{\rm pre}).
\]
For the second failure apply this after the first success and integrate
its exact record.  The two failure branches are disjoint, so
\[
 \Prob(J_k=1,A_{k+1}\mid\mathcal F_k)
       \le\varrho\Prob(J_k=1\mid\mathcal F_k),\qquad
 \Prob(J_k=0\mid\mathcal F_k)\ge c_0^2.                    \tag{G5b}
\]
The second bound uses at most two (G3) successes.  No hazard bound is
asserted after specifying the complete failed intermediate record.

For later division, consider any reachable nondead exact pair with radii
in \({\cal B}\), including a pair where \(p\ge a_0\) has just been
attained.  In a separate prospective continuation update the smaller
physical separation until both are at least \(p_*R\), where
\(p_*=1/200\).  Provisionally stop also at the first exit from the (G3)
input buffer \(r>3.3R\), \(s<1.7R\).  At every prospective update before
that exit, the selected quality is less than \(p_*R/\rho\le p_*=.005\).
Thus the cap is inactive and each success multiplies that record's physical
separation by at least \(1.8>\beta\).  Between its own updates its radius
and intrinsic quality stay fixed, since the all-success pair remains live.
For either record every prefix of its own starting-separation sequence
therefore has sum less than
\[
 {\beta\over\beta-1}p_*R=.015R.
\]
An elementary radius changes by exactly ten times that starting physical
separation.  Hence every prospective prefix has
\[
 r_{\rm f}>3.35R,\qquad s_{\rm f}<1.65R.                   \tag{G5c}
\]
This contradicts the proposed first exit from the (G3) input buffer, so
the continuation reaches its quality target within that buffer.  At the
end both normalized qualities
are at least \(p_*/4\), the ratio lies in \((9/8,5)\), and
\(6s_{\rm f}<9.9R<12R\).  Lemma~\ref{g:hookup} applies.  There are at
most \(2\max\{0,\lceil\log_\beta(p_*/p)\rceil\}\) successful updates.
With \(\alpha=2\log_\beta(c_0^{-1})\), this proves
\[
 \Prob(Z^\sigma\mid\mathcal F)
       \ge {c(p\wedge p_*)^\alpha\over\log R}.             \tag{G5}
\]
The prospective continuation is used only for this lower bound and may
leave \({\cal B}\); (G5c) supplies its larger buffer.

Set all flags after \(\tau\) to zero.  Freeze the state after a nondeath
stop, and keep its live indicator zero after death.  Let
\(N_k=\sum_{j<k}J_j\).  By (G5b),
\[
                  W_k=\mathbf1_{A_k}\varrho^{-N_k}
\]
is a nonnegative supermartingale.  For
\(I_n=(\beta^{-n},\beta^{1-n}]\), put
\[
 M_n=\sum_{k<\tau}\mathbf1_{\{A_k,\ p_k\in I_n\}},
\]
and let \(T_{n,m}\) be the \(m\)-th counted start.  Between two visits
to a band there is a completed flagged round: without a flag each
transition multiplies \(p\) by at least \(\beta\), so it cannot return.
The flag may occur in another band.  Optional stopping at bounded times,
followed by monotone convergence of the hitting indicators, gives on
\(S=T_{n,1}<\infty\)
\[
 \Prob(T_{n,m}<\infty\mid\mathcal F_S)\le\varrho^{m-1}.    \tag{G5d}
\]
Here \(\mathcal F_S\) is the stopped sigma-field on the first-visit event.
At a live start in \({\cal B}\), the ratio \(r/s\) lies in \([7/3,4]\)
and \(s\in[R,3R/2]\), so the upper bound in Lemma~\ref{g:hookup} is
\(C/\log R\).  Apply it at \(T_{n,m}\), use (G5d), and divide by (G5)
at the first visit, where \(\beta^{-n}<p<a_0<p_*\).  Integrating that first-visit
record yields
\[
 \Prob(M_n\ge m\mid\Phi_{4R},\Theta_R,Z^\sigma)
       \le C\beta^{\alpha n}\varrho^{m-1}.                 \tag{G6}
\]

Choose \(A>\alpha\log\beta/(-\log\varrho)\).  For some \(\xi<1\),
(G6) gives
\(\Prob(M_n>\lceil An\rceil\mid\Phi_{4R},\Theta_R,Z^\sigma)
\le C\xi^n\).  A failed physical separation satisfies
\(\ell_1=\rho_1q_1<2(1+10eq)\ell\le2.2\ell\).  The sums of elementary
physical inputs in the three possible round types are consequently less
than \(3.2\ell\), \(2.5\ell\), and
\(\ell+1.5\ell+2.2(1.5\ell)=5.8\ell\), respectively.  Thus the sum of
outer loss and inner gain in one round is less than \(58pR\).
For \(p<a_0\le10^{-3}\), a complete round begun in \({\cal B}\) has
intermediate radii
\[
 r>3.442R>3.25R,\qquad s<1.558R<1.61R,
\]
even when its endpoint triggers stopping.  These are the strong (G4)
buffers; the selected normalized qualities before their growth attempts
were already shown to be less than \(1/500\), as (G4) requires.  Its failed
output need not remain the smaller physical separation.

Low-quality rounds have \(n\ge N+1\).  Outside the exceptional events,
\[
 \sum_{k<\tau}p_k
 \le (A+1)\sum_{n\ge N+1}n\beta^{1-n}
 =3(A+1)(N+3)\beta^{-N}.
\]
Choose \(N\) so large that the resulting radial movement is less than
\(R/2\) and \(C\sum_{n\ge N+1}\xi^n<\eps/3\).  The radii cannot then
leave \({\cal B}\) while \(p<a_0\).  An infinite live
low-quality run would contain infinitely many flags, since after a last
flag repeated multiplication by \(\beta\) reaches \(a_0\).  The
supermartingale shows that survival through infinitely many flags has
probability zero.  Conditional on \(Z^\sigma\), therefore, \(p\ge a_0\)
is attained before buffer exit with probability at least \(1-\eps/3\).

On that attainment event condition on the exact transcript \(H\).  Since
its radii lie in \({\cal B}\),
\[
 q_{\rm o}\wedge q_{\rm i}\ge a_0/4,\qquad p\wedge p_*\ge a_0.
\]
Apply the one-record late-collapse argument to the outer record down to
\(3R\), keeping the attained inner record fixed; apply its outward version
to the inner record up to \(2R\), keeping the attained outer record fixed.
These are two estimates under the same \(H\), not estimates conditioned
on the other final record.  The argument in
\cite[proof of Lemma~3.15, equations (3.40)--(3.42)]{SSdiscrete} applies
with these fixed buffers.  A pair of close strand tips is screened at
dyadic circles about its midpoint by the single-sign circular-clearance
barriers proved above.  Each required breach costs a fixed factor, giving
\(C_{a_0}t^{\zeta_{\rm s}}\) for separation \(tR\).  For a walk--strand approach, expose
the walk only to its first \(tR\)-approach and use Beurling up to the next
fixed intermediate circle, giving \(C_{a_0}t^{\zeta_{\rm w}}\).  Leave the remaining
bridge unexposed in both cases.  The augmented upper bound in (G1) then
gives the joint-with-hookup bound
\(C_{a_0}(t^{\zeta_{\rm s}}+t^{\zeta_{\rm w}})/\log R\).  No (G3) call is
made in this late step.  The field-radius pairs for the augmented (G1)
factors are \((3R,s_H)\) and \((r_H,2R)\), with respective ratios in
\([2,3]\) and \([7/4,2]\).  Their clearance requirements are
\(6s_H\le9R\) and \(6(2R)=12R\); the intermediate walk prefix is covered
by the augmented upper bound.  Divide by (G5), which at the attained pair
is at least \(c a_0^\alpha/\log R\).  Choose \(0<t<.03\) so small that, under the
same \(H,Z^\sigma\), the union bound for the two late-collapse events is
less than \(2\eps/3\).  The bad events here mean that a defining physical
separation is at most \(tR\).  Either final normalized quality at most
\(t/3\) implies its corresponding bad event.  Taking
\(c(\eps)=t/3\), together with attainment, proves (G2).
\end{proof}

\begin{proposition}[Comparison of sampled configurations]
\label{g:sample-comparison}
There are \(C<\infty\) and \(R_0<\infty\) such that the following holds
for \(R\geq R_0\). Take two domains
containing \(\mathfrak B_{12R}\), with observers outside that disk,
and any two external records \(\phi,\phi'\) at radius \(4R\) having
positive hookup probability. For every event \(H\) determined by
the initially unoriented \(\Theta_R\), with the same specified geometric
incident-edge label in the two experiments and no conditioning on the
landing side,
\[
 C^{-1}\Prob'(H\mid\Phi'_{4R}=\phi',Z'^\sigma)
 \leq \Prob(H\mid\Phi_{4R}=\phi,Z^\sigma)
 \leq C\Prob'(H\mid\Phi'_{4R}=\phi',Z'^\sigma).
 \tag{G7}
\]
The constant is independent of \(R\) and of additional unused annuli
outside radius \(4R\).
\end{proposition}

\begin{proof}
Before hookup conditioning, the field observations in each
\(\Phi_{3R}\) are only the signs adjacent to its two external
arms, supported outside radius \(3R-O(1)\). The event specifying
an oriented \(\beta_{2R}^\sigma\) tests vertices inside
radius \(2R+O(1)\). Apply Lemma~\ref{g:strand-comparison} with
inner disk of radius \(9R/4\) and outer disk of radius \(5R/2\).
Both base laws have no observations in this outer disk, their graphs
agree there, and their external arms connect scale \(3R\) to the
original boundary outside scale \(6R\), providing the anchors.
Thus the inner oriented-strand laws are comparable by a fixed factor
under these different external records. Summing this comparison over
the two possible orientation lifts gives the same bound for the
unoriented inner record; no landing side is added to that record or
conditioned on. At this step no
\(\Theta\)-record has already been imposed in either base law:
the inner record is precisely the additional event being compared.
The reversed-walk law is also comparable by a
fixed factor: a walk from either entrance point has probability
bounded below to visit the other before reaching zero. Apply the
strong Markov property at that visit. The field and walk factors
are independent before hookup. Thus there is a reference weight
\(p_{2R}(\theta)\) for inner records such that their probabilities
given a record at radius \(3R\) are comparable to this weight.

Choose the quality threshold in Lemma~\ref{g:separation} with error
\(1/2\), and let \(X\) be the inner records at \(2R\) exceeding it.
Let \(\mathcal Q\) be the event that both the \(2R\) inner and
\(3R\) outer records exceed it. Separation holds even conditional
on a specified \(\Theta_R=\vartheta\). Hence conditioning further
on \(\mathcal Q\) changes the probability of this specified inner
record by at most a fixed factor. On \(\mathcal Q\), apply
Lemma~\ref{g:hookup} with inner radius \(2R\) and outer radius \(3R\).
Its ratio condition holds, and its clearance requirement is precisely
\(\mathfrak B_{12R}\), as assumed. Bayes' formula and the reference-weight
comparison give
\[
 \Prob(\Theta_{2R}=\theta\mid\mathcal Q,Z^\sigma,\Phi_{3R})
 \asymp \frac{p_{2R}(\theta)}{\sum_{\xi\in X}p_{2R}(\xi)}
 \quad(\theta\in X).
\]
Indeed the numerator before normalization is comparable to
\(p_{2R}(\theta)/\log R\), and summing determines its normalization.
Average over \(\Phi_{3R}\), then sum over all \(\theta\in X\)
whose restriction to radius \(R\) is \(\vartheta\).
The resulting reference ratio is independent of the external
experiment, proving (G7) first for single records, then for events.
Conditioning on data farther out and averaging over the radius
\(4R\) records preserves the same constant.
\end{proof}

\begin{lemma}[No additional sampled excursions]\label{g:extra}
Let \(\mathcal J(r,R)\) be the event that more than two disjoint
interface arcs cross from \(\mathfrak B_r\) to
\(\partial\mathfrak B_R\), or that the walk exits
\(\mathfrak B_R\) after first entering \(\mathfrak B_r\) and before
its landing at zero. For every \(\eps>0\) there are
\(A(\eps)<\infty\) and \(r_0(\eps)<\infty\)
such that
\[
 \Prob(\mathcal J(r,R)\mid\Phi_R,\Theta_r,Z^\sigma)<\eps
 \tag{G8}
\]
when \(r\geq r_0(\eps)\), \(R>A(\eps)r\), the observer and boundary are outside
\(\mathfrak B_{4R}\), and the conditioning has positive probability.
The same holds, with a harmless change of \(\eps\), conditioned
only on \(\Phi_R\) and \(Z=\{S_\tau=0\}\).
\end{lemma}

\begin{proof}
Apply Lemma~\ref{g:separation} in the outer buffer with base radius
\(R/4\); its clearance requirement is \(3R<4R\).  In the inner buffer
the base radius is \(r\), and its requirement \(12r<4R\) holds after
increasing \(A(\eps)\) so that \(A(\eps)>4\).  Separation in these
buffers makes all four records at scales \(2r,3r,R/2,3R/4\) have quality
at least \(c(\eps)\), except with probability \(O(\eps)\). Put \(\rho\asymp\sqrt{rR}\).
Let \(H=(\Phi_{3r},\Theta_{2r})\). The walk portion \(S_*\) from its
first radius-\(3r\) entrance to the last visit, before zero, to the
endpoint of its radius-\(r\) reversed piece is unlikely to leave
\(\mathfrak B_{\rho/3}\) without hitting the external arms. The
ordinary Beurling estimate and first-hit/last-exit decomposition
bound this probability, conditional on \(H\), by
\(b(R/r)\to0\). If it escapes, reveal only its prefix \(S_{**}\)
up to first exit. The augmented bound (GK), not a bound conditioned
on the completed walk, gives
\[
 \Prob(S_*\not\subset\mathfrak B_{\rho/3},Z^\sigma\mid H)
       \leq C b(R/r)/\log r.
\]
Divide by the lower bound \(c(\eps)/\log r\) on separated records
and average over them. This proves confinement under the landing law.

For the outer construction, use only \(\operatorname{ext}_{3r}S\).
Let \(S'\) be its terminal path from its last exit of
\(\mathfrak B_{R/3}\) to its endpoint at radius \(3r\).
Choose measurably from \(S'\) an arc on the circle of radius \(\rho\)
whose union with \(S'\) separates the inner and outer circles.
Intersect this open circular arc with \(D(\operatorname{ext}_{3r}\gamma)\).
Its two components next to \(S'\) have other endpoints on the boundary
of this deleted-triangle domain, not on the interface itself;
call them \(\alpha_1,\alpha_2\). Thus an endpoint attachment does
not already count as an interface intersection. They are
\(\Phi_{3r}\)-measurable. For each \(i=1,2\) separately, in the
component on that side place nested barriers separating \(\alpha_i\)
from the inner circle.
The deterministic construction is the clearance construction in
\cite[proof of (3.48)]{SSdiscrete}: take a separating circular arc;
if its middle has adequate boundary clearance, attach shortest
endpoint segments, and otherwise repeat at the smaller distance
between its two facing pure boundary pieces. The procedure
terminates because the two boundary pieces are disjoint.
Use the shortest endpoint segments and fixed relative endpoint
radius of the circular-clearance layout above. Thus each output
satisfies (UB1), although its lattice diameter may be small.
Barriers selected in distinct dyadic annuli remain separated.

For a fixed \(i\), after exposing the two interface arms to one
barrier, the next single-sign barrier satisfies
Proposition~\ref{g:barriers}, including its bounded-scale case.
Thus crossing \(N\) barriers in that access has probability at most
\((1-c)^N\), conditional on \(H\) but not on the landing.
Apply this separately to the two accesses, without conditioning on
the other shielding event; their union costs a factor of two.
For \(F=\{\gamma\cap(\alpha_1\cup\alpha_2)\ne\varnothing\}\),
fix the completed field while retaining the unexposed walk bridge.
The first augmented upper bound following (G1) gives
\[
 \Prob(F,Z^\sigma\mid H)
 =\E[1_F\Prob(Z^\sigma\mid H,\gamma)\mid H]
 \leq \frac{C}{\log r}\Prob(F\mid H).
\]
The same lower normalization gives \(C(\eps)(1-c)^N\) under the
landing law, where \(N\to\infty\) with \(R/r\). Repeat the whole
construction from the inner reversed record \(\Theta_{R/3}\),
without also revealing the forward intermediate bridge.

Outside these exceptional events, confinement identifies the terminal
forward path \(S'\) with the corresponding portion of the reversed
inner record. The same deterministic circular-arc selection then
makes the two crosscuts from the separate descriptions coincide. Their
boundary endpoints therefore attach to lattice edges incident to both
the inner incident strand and the two external arms. These strand
attachments are on opposite sides of the incident
edge: the other possibility contradicts the winding number of
the circular separator completed along the common walk path.
This is the deterministic Jordan-curve conclusion in
\cite[final paragraphs of Lemma 3.17]{SSdiscrete}; its hypotheses
are the two coinciding crosscuts, a simple interface, and the
common walk avoiding it before zero, all obtained above.
Thus each external arm has merged with a distinct inner arm,
and there can be no additional annular interface crossing.

The total error is \(C\eps+o_{R/r}(1)+C(\eps)(1-c)^N\).
Choose the buffer error first and then \(R/r\) large.
Average over \(\Theta_r\) and the six possible incident edges
for the external-only statement.
\end{proof}

\subsection{Boundary clearance, survival, and nonreturn}

The survival and nonreturn arguments adapt
\cite[Lemma~3.26, especially equation (3.62)]{SSdiscrete}, using the
barrier and sampled-excursion estimates proved above.
For the boundary-attaching arguments in this subsection, fix
\(\lambda_*>0\) and assume the exterior Dobrushin data satisfy
\[
 h_\partial\in[\lambda_*,1/2]\text{ on }\partial_+,
 \qquad
 h_\partial\in[-1/2,-\lambda_*]\text{ on }\partial_-.
 \tag{GB}
\]
All constants may depend on \(\lambda_*\). This includes the reference
data \(\pm1/2\), and the candidate matched data \(\pm l\) with
\(l\geq c\) from the completed-reference use of Lemma~\ref{g:sign-trace}
in Proposition~\ref{h:height-gap}. The boundary-landing argument itself
uses the prescribed exterior lower bound in \((\mathrm{GB})\), not a
new application of the string comparison on an open sector. No uniform assertion as
\(l\downarrow0\) is required. The source's boundary condition
\cite[condition \((\partial)\), Section 3.7]{SSdiscrete} is more
permissive, allowing small wrong-sign heights; we do not need or invoke
that additional Gaussian statement.

\begin{lemma}[Clearance from the original boundary]
\label{g:boundary-landing}
Let \(S\) be an independent full walk from \(v\), stopped only on
the original boundary at \(\tau_D\), and let \(\gamma\) be the
completed interface. Put \(d=\operatorname{dist}(v,\partial D)\) and
\[
 M=\{S_j:j<\tau_D,\ S_j\text{ is adjacent to }\gamma\}.
\]
Under \((\mathrm{GB})\), uniformly over domains, observers, and the
indicated exterior data,
\[
 \lim_{\eta\downarrow0}\limsup_{d\to\infty}
 \Prob\bigl(\operatorname{dist}(M,\partial D)<\eta d\bigr)=0.
 \tag{GB1}
\]
The distance of the empty set is infinity. In particular the assertion
holds for a terminal harmonic landing, or simultaneously for every
old-prefix landing of this walk that is not on \(\partial D\).
\end{lemma}

\begin{proof}
We adapt the full-walk argument of \cite[Lemma~3.23, especially
equations (3.53) and (3.55)]{SSdiscrete}. We first establish pure-arc
repulsion. For \(y\in\partial_+\), choose
any deterministic lower bound \(b\leq d_D(y,\partial_-)\).
Start the fixed-ratio iteration at \(\rho\le b/8\), losing only
\(O(1)\) outer scales. Take the component of \(B(y,\rho)\cap D\)
accessible from \(y\), and its circular crosscut separating \(y\)
from \(\partial_-\). Use the clearance construction of
Lemma~\ref{g:extra}, with \(\delta\le1/100\) and
\(c_A\le\delta/100\) in the circular-clearance layout. Write
\(\beta_1,\beta_2\) for the two components of
\(\partial_+\setminus\{y\}\). Its endpoint \(z_j\) is accessible
from \(y\) by a path of diameter at most
\(2(1+2\delta)\rho+O(1)\). Its internal separation \(\sigma_j\)
from \(\beta_{3-j}\) is greater than \(\ell_{\rm att}\) in the
wide case, and in the narrow case is at least
\(\delta\,\delta^k\rho\), with
\(\ell_{\rm att}\le\delta^k\rho\). For the endpoint radius
\(d_{\rm e}=c_A\ell_{\rm att}\le\delta\rho/100\), in both cases
\(\sigma_j>3d_{\rm e}\).

Let \(\chi_j\) be the first circular crosscut at radius \(d_{\rm e}\)
met by the initial segment from \(z_j\). Each endpoint of \(\chi_j\)
joins \(z_j\) along that segment and \(\chi_j\) by a route of
diameter at most \(2d_{\rm e}\). It cannot lie on \(\beta_{3-j}\), by
\(\sigma_j>3d_{\rm e}\). It cannot lie on \(\partial_-\), since concatenating
with the preceding access from \(y\) gives diameter at most
\[
       2(1+2\delta)\rho+2d_{\rm e}+O(1)<b
\]
at diverging scales, contrary to \(b\le d_D(y,\partial_-)\).
Prior exposed arms stay outside by the annular separation.
Thus both endpoints of \(\chi_j\) lie on \(\beta_j\).
The barrier tail avoids \(\chi_j\): its own nearest segment crosses
the circle once, and the wide middle has clearance
\(\ell_{\rm att}>d_{\rm e}\). If the other nearest segment met \(\chi_j\),
its remaining length to \(\beta_{3-j}\) would be at most \(d_{\rm e}\);
together with the preceding \(2d_{\rm e}\) route this would contradict
\(\sigma_j>3d_{\rm e}\). Hence \(\chi_j\) and the intervening
\(\beta_j\) subarc cut off the pure side containing the initial
prefix. This is its endpoint pocket \(A_{z_j}\).
At each stage this is one positive barrier, so no mixed layout is used.
Explore the interface only up to the next outer crosscut. Its exact
sign transcript leaves the inner barrier compatible, so
Proposition~\ref{g:barriers} gives a fixed conditional chance of
avoiding this crosscut and every smaller access to \(y\).
Iterating, then reflecting for \(y\in\partial_-\), gives constants
\(C,\zeta>0\) such that
\[
 \Prob\{d_D(y,\gamma)<r\}
 \leq C(r/b)^\zeta+o_d(1),                         \tag{GB2}
\]
when \(r,b\) are fixed positive multiples of \(d\) and \(r<b\).
The lattice error is uniform for these fixed multiples. This is the
deterministic crosscut argument of
\cite[Theorem 3.22]{SSdiscrete}, with its field input replaced by
our exterior version of Proposition~\ref{g:barriers}.

Fix \(r=\eta d\), and condition on \(\gamma\). At the first walk
visit to \(M\) at distance less than \(r\) from \(\partial D\),
the strong Markov property and the planar walk estimate give a
fixed probability of exiting \(D\) in the ball of radius \(2r+O(1)\)
about a closest boundary point. One can force the walk to surround
that point inside this ball, which forces the exit. Its continuation
and the adjacent interface edge then give an internal path of diameter
at most \(4r+O(1)\) from the exit \(Y=S_{\tau_D}\) to \(\gamma\).
Consequently
\[
 \Prob\{\operatorname{dist}(M,\partial D)<r\}
 \leq C\Prob\{d_D(Y,\gamma)<4r+O(1)\}.             \tag{GB3}
\]
Here \(Y\), unlike the interface landing, is independent of
\(\gamma\). Conditional on \(Y\), apply (GB2) whenever its internal
distance to the opposite arc exceeds \(s=\sqrt\eta\,d\), using
the deterministic lower bound \(b=s\) in that estimate. This gives
\(C\eta^{\zeta/2}+o_d(1)\).

It remains to bound exits within internal distance \(s\) of both
sign arcs. Delete all paths of diameter at most \(2s\) joining the
two arcs and retain the component of \(v\). Its new boundary lies
in two balls of radius \(4s\). To verify this elementary fact,
classify each deleted crosscut by which marked sign-change endpoint
it separates from \(v\). Any two crosscuts in the same class which
touch the retained component intersect: otherwise their separation
order prevents one of them from touching that component. Their
diameter bound puts the whole class in one ball. There are two
classes. An exit near both arcs is beyond this new boundary.
The walk hitting estimate for these two balls gives
\(C(s/d)^{1/2}=C\eta^{1/4}\). Combining with (GB3) proves (GB1).

Every nonexterior old-prefix landing is a vertex of \(M\), pathwise,
because the prefix is contained in \(\gamma\) and the full walk has
not yet exited \(D\). This proves the simultaneous assertion without
conditioning on any landing event or on any stopping time.
\end{proof}

\begin{lemma}[Clearance from the observer]\label{g:observer-landing}
Uniformly over the same bounded-data experiments,
\[
 \lim_{\eta\downarrow0}\limsup_{d\to\infty}
 \Prob\{\operatorname{dist}(v,\gamma)<\eta d\}=0.     \tag{GB4}
\]
This assertion does not require the lower bound in \((\mathrm{GB})\).
\end{lemma}

\begin{proof}
The nested shielding argument follows \cite[Lemma~3.24]{SSdiscrete},
using Proposition~\ref{g:barriers}. Expose the first interface approach
to the circles centered at \(v\) of radii \(2^{-j-1}d\).
After the first approach to radius \(\rho\), use the free closed
circle of radius \(3\rho/4\) as one single-sign barrier.
It has no endpoints, and has fixed relative clearance from both
the observed strands and the original boundary. Thus (UB1) holds
without an attachment module. Proposition~\ref{g:barriers} gives a
fixed chance of protecting the next smaller disk. Failure through \(N\) scales
has probability at most \((1-p)^N\). Fix \(N\) first and let
\(d\to\infty\), so every used scale diverges; then let \(N\to\infty\).
This proves (GB4). It also excludes all terminal and old-prefix
landings within \(\eta d\) of \(v\), up to one lattice spacing.
\end{proof}

\begin{lemma}[Survival of a stopped harmonic landing]\label{g:survival}
Let \(T\) be any stopping time for the filtration generated by oriented
initial interface pieces, \(K=\gamma[0,T]\), and
\(d=\operatorname{dist}(v,\partial D)\). Let \(X_T\) be the first
landing of an independent full walk \(S\) on \(\partial D(K)\);
let \(Y\) be its landing for the completed interface. For every
\(\eps>0\) there are \(p>0\), \(d_0<\infty\) such that
\[
 \Prob\bigl(\Prob(Y=X_T\mid K,S)<p\bigr)<\eps
 \qquad(d>d_0),                                      \tag{G9}
\]
uniformly over the domains, observers, and stopping times under
\((\mathrm{GB})\).
\end{lemma}

\begin{proof}
Conditional on an exact stopped prefix \(K\), the field has the
ordinary strand-conditioned law: whether the rule stopped at this
prefix is already decided by that prefix. The independent walk
adds no field conditioning. Put \(U=D(K)\).
For small fixed \(\delta>0\), except on an event tending uniformly
to zero with \(\delta\), the following four conditions hold:
\[
\begin{split}
 &d_U(X_T,\hbox{opposite-sign side})\geq2\delta d,\\
 &\operatorname{diam} S[0,\tau_T]<\delta^{-1}d,\\
 &\operatorname{diam} S[t_*,\tau_T]<\delta d/4,
 \quad t_*=\inf\{t:d(S_t,\partial U)\leq\delta^2d\},\\
 &d(v,K)\geq\delta^{1/2}d.
\end{split}                                                   \tag{G10}
\]
The middle two statements are the planar walk hitting estimate
uniform in \(U\). The last statement follows from
Lemma~\ref{g:observer-landing} and \(K\subset\gamma\).
For the observer-separated stopping times needed in the driving
argument, the last condition can instead be imposed deterministically.

For the first condition delete all crosscuts of diameter at most
\(4\delta d\) joining the two sign sides. The new boundary of the
component containing \(v\) lies in two balls of radius \(8\delta d\).
Indeed each deleted crosscut separates \(v\) from one of the two
sign-change points; any two crosscuts for the same point which
touch the retained component must intersect. All such crosscuts
are consequently within twice their diameter of a fixed one.
Apply the walk hitting estimate to those two balls, using the last
condition. This also explains why arbitrary narrowing of \(U\)
does not invalidate (G10).

Fix records satisfying (G10), and suppose the landing is positive.
The deterministic clearance construction used in
Lemma~\ref{g:extra} gives a compatible crosscut \(\Upsilon_0\)
in \(B(X_T,\delta d)\setminus B(X_T,\delta d/2)\), joining the
positive boundary pieces on either side of \(X_T\) and separating
it from the negative side. If no hexagon centered at a vertex of the
stopped walk meets this crosscut,
avoidance of this crosscut protects the entire walk.
Otherwise add the boundaries of the walk hexagons on its
unprotected side and take the boundary facing the negative side.
Its diameter is at most \(2\delta^{-1}d\); its new part is at
distance at least \(\delta^2d/2\) from \(\partial U\), by the
terminal-walk condition. Indeed a contacting walk hexagon occurs
before \(t_*\), since the terminal walk piece lies in
\(B(X_T,\delta d/4)\), whereas the crosscut lies outside
\(B(X_T,\delta d/2)\). That contact is at distance at least \(\delta^2d-O(1)\)
from \(\partial U\); therefore the boundary-attached original
crosscut has diameter at least \(\delta^2d-O(1)\) as well.
The old endpoints retain their pure pockets and their shortest
initial segments. In the contacting case choose their endpoint
radius \(d_{\rm e}\) also below \(\delta^2d/(100A)\).
The fixed relative circular construction and the preceding diameter
bound make this a fixed further shrink depending only on \(\delta\).
The added part cannot enter \(B(z,d_{\rm e})\) and has clearance
much larger than \(d_{\rm e}\). Together with the circular-clearance
qualification this proves (UB1) for the modified crosscut.
It is one landing-sign barrier with parameter \(c\delta^3\)
after a fixed shrink, at scale \(2\delta^{-1}d\).

Proposition~\ref{g:barriers} gives probability at least \(p_\delta>0\)
that the future interface avoids the resulting barrier. On this
event it is adjacent to no vertex of the stopped walk, so \(Y=X_T\).
In the first branch \(\Upsilon_0\) can have bounded lattice diameter;
this is covered by the all-scale part of that proposition.
Reflect signs for a negative landing. Choose \(\delta\) to make
the failure probability of (G10) smaller than \(\eps\), and then
take \(d\) sufficiently large.
\end{proof}

\begin{proposition}[Nonreturn at stopped harmonic landings]
\label{g:nonreturn}
For every \(\eps>0\) there are \(s>0\) and \(d_0<\infty\) such that
\[
 \Prob\bigl(d(X_T,\gamma\setminus\gamma[0,T])<s d,\,
                  X_T\notin\partial D\bigr)<\eps
 \qquad(d>d_0),                                      \tag{G11}
\]
uniformly over all sign-filtration stopping times \(T\), with
\(d=\operatorname{dist}(v,\partial D)\), under \((\mathrm{GB})\).
\end{proposition}

\begin{proof}
Let \(\mathcal G=\sigma(K,S)\), \(A=\{Y=X_T\}\), and
\(q=\Prob(A\mid\mathcal G)\). Lemma~\ref{g:survival} implies,
for every \(\mathcal G\)-measurable event \(E\),
\[
 \Prob(E)\leq\eps_1+p^{-1}\Prob(E\cap A).              \tag{G12}
\]
Take \(E\) to be the event that the old prefix has more than two
crossings of an annulus about its walk landing. On \(A\), this
implies the corresponding event for the completed interface at
its completed harmonic landing. Lemma~\ref{g:extra}, with error
smaller than \(p\eps_1\), therefore removes all additional old
excursions from a small comparison window.

Exclude landings close to the old tip using the walk hitting bound.
Lemmas~\ref{g:boundary-landing} and \ref{g:observer-landing} exclude
landings close to the original boundary or observer, uniformly over
all prefixes. Choose the largest comparison window smaller than
one tenth of this retained clearance. This verifies the original
boundary and observer exclusions when Lemma~\ref{g:extra} is applied
after recentering at the landing.

In the remaining window exactly two incident old arms reach the
outer circle, and the old tip lies outside it. There are two
relevant complementary components. For each component separately,
under the same conditioning \((K,S)\), select nested single-sign
circular-clearance barriers attached to its pure old side, with
the shortest endpoint segments and fixed relative endpoint radii
specified above. After each failed shield, reveal the future arms
up to that barrier; the conditional field law is again an
exact-strand law, so the next shield succeeds with probability at
least a fixed \(c>0\). Failure through \(N\) scales in that access
has probability at most \((1-c)^N\). We do not condition one access
estimate on the other shielding event; the union bound costs two.
When neither complementary component is traversed, the future
cannot enter the innermost Euclidean ball. The removal of
additional old excursions is what converts this internal-access
conclusion to Euclidean nonreturn.

Choose \(\eps_1\), obtain \(p\) from survival, choose the completed
excursion error below \(p\eps_1\), choose \(N\), and then let all
the finitely many retained scales tend to infinity in lattice
units. Finally choose their innermost relative radius \(s\).
The sum of the discarded probabilities can be made less than
\(\eps\), proving (G11).
\end{proof}

\begin{remark}[What the geometric result supplies]
The random mass of old-boundary accesses revisited by the future is
the conditional harmonic probability of the event in (G11).
Its expectation tends to zero, uniformly over stopping rules.
Consequently bounded harmonic boundary tests may discard this mass.
This is the estimate needed before a tower-property argument for a
stopped observable. Absolute continuity of access in each fixed domain
alone would not give its uniform form.
\end{remark}
